% ============================================================================
%  FINAL SUBMISSION VERSION  —  Phase 2.2 (V3), Introduction revised (Phase 3)
%  Title : Joint Code-Dimension and k-Galois-Hull Enumerators for Simple-Root
%          Constacyclic Codes over Square-Free Affine Algebras
%  Author: Arun Kumar Yadav, K.R. Mangalam University, Gurugram, Haryana, India
% ----------------------------------------------------------------------------
%  Files required to compile:  references.bib, sn-mathphys-num-local.bst
%  Build order:  pdflatex FILE  ->  bibtex FILE  ->  pdflatex FILE  ->  pdflatex FILE
%  Result of the Phase-2.2 build: 26 pages, 0 errors, 0 overfull boxes,
%  0 undefined references. That build predates the Phase-3 revision of the
%  Introduction; the Introduction was rewritten (same mathematics, same
%  citation set) and recompiled must be re-checked for page count and overfull
%  boxes, because FINAL_SUBMISSION.bbl was reordered to match the new order of
%  first citation. All figures are inlined vector TiKZ; no external images.
% ----------------------------------------------------------------------------
%  BEFORE UPLOAD (author action, tracked in SCIENTIFIC_ISSUES_TO_REVIEW.md):
%   * add the corresponding-author e-mail address (item D3);
%   * confirm the 2026 bibliographic fields at item C2, and the theorem number
%     of the entry debnathAffine2026 at item C1 against the published
%     Discrete Mathematics version;
%   * attach programs/ and computed_data/ as supplementary material, or adjust
%     the data-availability wording (item D2).
% ============================================================================
\documentclass[10pt,a4paper]{article}

\usepackage{iftex}
\ifPDFTeX
  \usepackage[T1]{fontenc}
  \usepackage[utf8]{inputenc}
\fi
\usepackage{lmodern}
\usepackage{amsmath,amssymb,amsthm,mathtools}
\usepackage{mathrsfs}
\usepackage{booktabs}
\usepackage{array}
\usepackage{tabularx}
\usepackage{geometry}
\usepackage[numbers,square,sort&compress]{natbib}
\usepackage{hyperref}
\providecommand{\bibcommenthead}{}
\usepackage{microtype}
\usepackage{tikz}
\usetikzlibrary{arrows.meta,positioning}
\usepackage{caption}
\captionsetup[figure]{name=Fig.,labelfont=bf,labelsep=space}
\usepackage{titlesec}
\usepackage{enumitem}
\titleformat{\section}{\normalfont\large\bfseries}{\thesection.}{0.6em}{}
\titleformat{\subsection}{\normalfont\normalsize\bfseries}{\thesubsection}{0.6em}{}
\titlespacing*{\section}{0pt}{13pt plus 2pt minus 2pt}{6pt plus 1pt}
\titlespacing*{\subsection}{0pt}{10pt plus 2pt minus 2pt}{4pt plus 1pt}
\setlist{topsep=4pt,itemsep=2pt,parsep=0pt}
\setlength{\bibsep}{0pt plus 1pt}
\renewcommand{\bibfont}{\small}
\usepackage{etoolbox}
\appto\normalsize{%
  \setlength{\baselineskip}{12.2pt}%
  \setlength{\abovedisplayskip}{5pt plus 2pt minus 1pt}%
  \setlength{\belowdisplayskip}{5pt plus 2pt minus 1pt}%
  \setlength{\abovedisplayshortskip}{3pt plus 1pt}%
  \setlength{\belowdisplayshortskip}{4pt plus 1pt minus 1pt}%
}
\urlstyle{same}
\geometry{textwidth=150mm,textheight=213mm,centering}
\hypersetup{colorlinks=true,linkcolor=blue,citecolor=blue,urlcolor=blue,
  pdftitle={Joint Code-Dimension and k-Galois-Hull Enumerators for Simple-Root Constacyclic Codes over Square-Free Affine Algebras},
  pdfauthor={Arun Kumar Yadav},
  pdfsubject={Coding theory}}

% Section numbers are followed by a period in headings, as in 1. Introduction.
\makeatletter
\renewcommand{\@seccntformat}[1]{\csname the#1\endcsname.\quad}
\makeatother

% Ragged-right paragraph columns, so that narrow table cells are not stretched.
\newcolumntype{L}[1]{>{\raggedright\arraybackslash}p{#1}}

\newcommand{\F}{\mathbb F}
\newcommand{\K}{\mathbb K}
\newcommand{\A}{\mathcal A}
\newcommand{\Fs}{\mathcal F_s}
\newcommand{\Os}{\mathcal O_s}
\newcommand{\Epoly}{\mathscr E}
\newcommand{\Hull}{\operatorname{Hull}}
\newcommand{\rank}{\operatorname{rank}}
\newcommand{\wt}{\operatorname{wt}}
\newcommand{\code}{\mathrm{code}}
\newcommand{\cand}{\mathrm{candidate}}
\newcommand{\perpE}{\perp_E}
\newcommand{\Dk}{D_k}

% Keep theorem separation visible without the article class's large list gaps.
\makeatletter
\def\thm@space@setup{%
  \thm@preskip=4pt plus 1pt minus 1pt%
  \thm@postskip=4pt plus 1pt minus 1pt}
\patchcmd{\proof}{\topsep6\p@\@plus6\p@\relax}
  {\topsep4\p@\@plus2\p@\relax}{}{}
\makeatother
\setlength{\emergencystretch}{1em}

% Numbered environments share one counter per section; equations are numbered
% by section as well, so that every reference carries its own location.
\numberwithin{equation}{section}
\theoremstyle{plain}
\newtheorem{theorem}{Theorem}[section]
\newtheorem{lemma}[theorem]{Lemma}
\newtheorem{proposition}[theorem]{Proposition}
\newtheorem{corollary}[theorem]{Corollary}
\theoremstyle{definition}
\newtheorem{definition}[theorem]{Definition}
\newtheorem{example}[theorem]{Example}
\theoremstyle{remark}
\newtheorem{remark}[theorem]{Remark}

% Thin rules mark the main theorem of the paper; used once, in Section 4.
\newenvironment{keyresult}
  {\par\addvspace{3pt}\noindent\rule{\linewidth}{0.4pt}\par\nobreak\addvspace{3pt}}
  {\par\nobreak\addvspace{3pt}\noindent\rule{\linewidth}{0.4pt}\par\addvspace{3pt}}

\title{Joint Code-Dimension and \texorpdfstring{$k$}{k}-Galois-Hull Enumerators for Simple-Root Constacyclic Codes over Square-Free Affine Algebras}
\author{\normalsize Arun Kumar Yadav\\
\small School of Basic and Applied Sciences,\\
\small K.R. Mangalam University, Gurugram, Haryana, India}
\date{}

\begin{document}
\maketitle


\begin{center}
\noindent\textbf{Abstract}
\end{center}
For compatible twists, this paper counts the simple-root constacyclic codes over a fixed square-free affine algebra $\A$ defined over a finite field, classified jointly by code dimension and $k$-Galois-hull dimension. Every such code is a product of component ideals determined by a labeled choice of irreducible factors of the component moduli. Under the compatibility condition, a Frobenius-twisted reciprocal permutes these factors, and the hull dimension becomes a weighted count of $1$-to-$0$ transitions in the selection word read along each orbit of this permutation. A two-state weighted transfer matrix, or the equivalent closed one-orbit polynomial, turns the orbit data into a bivariate generating polynomial whose coefficients are exactly the numbers of codes realizing each pair of code dimension and hull dimension. The polynomial yields the code-dimension and hull-dimension distributions, the exact count of complementary-dual codes, and the mean, variance and conditional distribution of the hull dimension at each attainable code dimension. The formulas agree with direct finite-field linear-algebra computations in characteristics $2$, $3$ and $5$ and recover published benchmark counts, and for $\F_{32}$ with $n=31$ they give the exact enumerator over all $2^{31}$ labeled factor selections without enumerating them.

\medskip
\noindent\textbf{Keywords:} $k$-Galois hulls; constacyclic codes; square-free affine algebras; exact labeled enumeration; generating polynomials; transfer matrices.

\medskip
\noindent\textbf{2020 MSC:} 94B05, 94B15, 05A15.

\section{Introduction}\label{sec:intro}

The hull of a linear code $C$ is the intersection $C\cap C^{\perp}$ with its
dual, and its dimension is a structural invariant of the code. Replacing the
Euclidean inner product by a Frobenius-twisted pairing in one argument leads to
the $k$-Galois inner product and to the $k$-Galois hull, of which the Euclidean
and Hermitian hulls are the cases $k=0$ and $k$ equal to half the degree of the
field \cite{fanZhang2017}. For cyclic-type codes the ambient space is a
polynomial quotient, so duals, intersections and hence hull dimensions are
accessible to factorization arguments \cite{huffmanPless2003}. Hull dimensions
therefore describe a single code and, at the same time, define a statistic of an
entire family; this paper studies the joint distribution of that statistic and
the code dimension.

The hull dimension, and not only its vanishing, controls several constructions.
The algorithms that test permutation equivalence of two linear codes, and those
that compute the automorphism group of a code, become expensive when the hull is
large, so the hull dimension measures their complexity \cite{sendrier1997}.
Codes with a trivial hull are the linear codes with complementary duals studied
by Massey \cite{massey1992}, and they can attain the Gilbert--Varshamov bound
\cite{sendrier2004}. In quantum coding the hull
dimension of a classical code determines the number of entangled pairs needed by
the corresponding entanglement-assisted quantum error-correcting code
\cite{brun2006,wildeBrun2008}, and the same mechanism is still used to obtain
quantum codes from constacyclic codes over rings \cite{zhang2026}.

The hulls of cyclic and negacyclic codes over finite fields were the first to be
described through the factorization of the modulus. Yang and Massey
characterized when a cyclic code has a complementary dual \cite{yangMassey1994},
Skersys computed the average dimension of the Euclidean hull over the cyclic
codes of a fixed length \cite{skersys2003}, and Sangwisut et al.\ determined the
Euclidean and Hermitian hull dimensions of cyclic and negacyclic codes in closed
form and counted the codes realizing each prescribed hull dimension
\cite{sangwisut2015}. Constacyclic codes, which contain both families, were
studied in the same manner: Jitman and Sangwisut obtained the average dimension
of the Hermitian hull of constacyclic codes over fields of square order
\cite{jitmanSangwisut2018}, and Debnath et al.\ extended the dimension formulas
to the $k$-Galois pairing for constacyclic codes over finite fields, together
with counts of the codes of a prescribed Galois hull dimension under
restrictions on the field parameters \cite{debnathConstacyclic2023}; a correction
appeared in \cite{debnathCorrection2023}. This line continues with the average
Galois hull dimension of constacyclic codes \cite{debnathAverage2025} and with
the small values a Galois hull dimension can attain \cite{debnathSmall2026}.

The Galois pairing also made the hull accessible for codes that are not
constacyclic. Liu and Pan characterized the Galois hull of an arbitrary linear
code over a finite field \cite{liuPan2020}, Cao constructed MDS codes with
Galois hulls of prescribed dimensions \cite{cao2021}, and Fu and Liu treated
Galois self-orthogonal constacyclic codes \cite{fuLiu2022}. More recently,
generator descriptions have been derived that avoid passing to the dual code:
generator polynomial matrices for the Galois hulls of multi-twisted codes
\cite{eldinSole2026}, and generator matrices for the intersection of a linear
code with its Galois dual \cite{eldinLeroy2026}, the operation that defines the
hull.

A separate question is not what the hull of one chosen code is, but how hull
dimensions are distributed over a family. For unrestricted linear codes over a
finite field, a mass formula gives the number of codes of a fixed length and
Euclidean hull dimension \cite{sendrier1997}, and the sequences of these counts
as the hull dimension grows have been studied recently \cite{bouyuklieva2026}.
Distributions and asymptotics have also been obtained for double cyclic codes
\cite{gao2023} and for double circulant and four circulant codes with small hull
dimension \cite{aliabadi2026}. The averages computed in \cite{skersys2003},
\cite{jitmanSangwisut2018} and \cite{debnathAverage2025} are the first moments of
such distributions and do not determine them.

When the alphabet is a ring rather than a field, the factorization theory
changes. The structure theory of cyclic and negacyclic codes over finite chain
rings \cite{dinhLopezPermouth2004} and the Hamming-distance theory over such
rings \cite{nortonSalagean2000} provide the background for the analyses below.
Talbi et al.\ studied the hulls of cyclic serial codes over a finite chain ring
\cite{talbi2022}; Jitman et al.\ treated cyclic codes over $\mathbb{Z}_4$, where
the length need not be coprime to the characteristic \cite{jitmanZ4}; and Pathak
and Sharma covered the oddly even case, which requires separate arguments
\cite{pathakZ4}. Constacyclic codes over a finite non-chain ring, with
applications to quantum codes, were considered in \cite{zhang2026}. Closest to
the present setting, Debnath et al.\ studied $k$-Galois hulls of constacyclic
codes over affine algebra rings, exhibiting generators of the duals and hulls, a
hull-dimension formula, LCD criteria and quantum-code applications
\cite{debnathAffine2026}.

Two kinds of results have therefore been obtained: descriptions of the hull of
an individual code, and counts or averages of a single statistic of a family.
Neither kind determines the joint distribution of the code dimension and the
hull dimension, and it is convenient to separate three questions: (i) which hull
dimensions occur; (ii) how many codes realize a prescribed hull dimension; and
(iii) how many codes realize each pair $(\kappa,\eta)$ formed by a code
dimension $\kappa$ and a hull dimension $\eta$. Questions (i) and (ii) are
answered for cyclic and negacyclic codes in \cite{sangwisut2015} and, over
finite fields, for constacyclic codes in \cite{debnathConstacyclic2023}, where
the counts require restrictions on the field parameters, while for unrestricted
linear codes with respect to the Euclidean dual, question (iii) is answered by
\cite{sendrier1997} and \cite{bouyuklieva2026}. Over an affine algebra only
question (i) has been answered, by the hull-dimension formula of
\cite{debnathAffine2026}. Such a formula does not record the code dimensions at
which the hull dimensions occur, and a hull-only count cannot be aggregated over
the different factor selections into a joint distribution. Answering question
(iii) in this setting therefore requires carrying the code dimension along with
the hull dimension and treating cycles of the factor permutation that are longer
than the fixed points and two-cycles that arise in the Euclidean and Hermitian
cases.

Motivated by this gap, we determine in this paper the joint distribution of code
dimension and $k$-Galois hull dimension for simple-root constacyclic codes over
a fixed square-free affine algebra $\A$ with labeled field components. The
counted objects are labeled factor selections, in bijection with the
$\lambda$-constacyclic codes over $\A$ (Lemmas~\ref{lem:complete}
and~\ref{lem:global-code}); each is a distinct submodule of the fixed ambient
space, and no isometry classes are formed. The output is a bivariate generating
polynomial in which the exponents of $u$ and $z$ record the code dimension
$\kappa$ and the $k$-Galois hull dimension $\eta_k$, so that its coefficients
answer the three questions above simultaneously. The method is to
decompose $\A$ into its field components, to translate $\lambda$-constacyclic
codes into labeled selections of irreducible factors of the component moduli,
and to follow the $k$-Galois dual through the Frobenius-twisted reciprocal
permutation of those factors. Both statistics are additive over the orbits of
this permutation, and this is what makes the enumeration factorize.

The results of the paper are the following.
(i) The $k$-Galois dual of a component code is generated by the normalized
inverse-Frobenius reciprocal of the parity-check polynomial and is constacyclic
with twist $\lambda_s^{-p^{d_s-k}}$ (Theorem~\ref{thm:dual-generator}), obtained
from the classical Euclidean description of a constacyclic dual
(Lemma~\ref{lem:ordinary-reciprocal}) by the coordinatewise inverse Frobenius
(Proposition~\ref{prop:dual-translation}).
(ii) Under the compatibility condition \eqref{eq:compat} the reciprocal
permutes the irreducible factors of the component modulus. The hull support of a
code is then a cyclic boundary statistic on the orbits of that permutation, and
the hull dimension is a weighted count of the resulting $1$-to-$0$ transitions
(Theorem~\ref{thm:hull-support} and Proposition~\ref{prop:boundary}).
(iii) Each orbit contributes a two-state weighted transfer matrix whose trace,
the orbit polynomial, also has a closed binomial form
(Propositions~\ref{prop:transfer-matrix}--\ref{prop:closed}).
(iv) The exact labeled joint $(\kappa,\eta)$ enumerator over the CRT product,
with orbit weights $w_O=m_sd_O$, is
\[
\Epoly(u,z)=\sum_{C}u^{\kappa(C)}z^{\eta_k(C)}
 =\prod_{s}\prod_{O\in\Os}\operatorname{tr}\bigl(T_{w_O}(u,z)^{a_O}\bigr)
\]
(Theorem~\ref{thm:central}), where the sum runs over all $\lambda$-constacyclic
codes over $\A$. Items (i)--(iii) assemble established descriptions of component
duals and closed-walk counting; the joint enumeration (iv) combines them into
one product formula. From it the two marginal distributions, the exact LCD
count, and the mean, variance and conditional distribution of the hull dimension
follow by specialization
(Corollaries~\ref{cor:total}--\ref{cor:conditional-variance}), while the
extension-degree weights $w_O$ convert component dimensions into dimensions over
the common base field.

The enumeration is exact: the coefficients of $\Epoly(u,z)$ are integers, the
conditional moments are reduced rationals, and for $q=32$, $n=31$ the orbit
recurrence evaluates all $2^{31}$ labeled factor selections without enumerating
them. For the families in characteristics $2$, $3$ and $5$ listed in
Section~\ref{sec:verification}, every $k$-Galois hull dimension computed
directly by finite-field linear algebra agrees with the orbit formula, and two
published benchmark counts \cite{sangwisut2015} are recovered. The orbit product
assumes simple-root moduli and compatible twists, and codes are counted as
labeled selections rather than isometry classes; the scope is stated at the end
of Section~\ref{sec:conclusion}. Figure~\ref{fig:pipeline} summarizes the
workflow. Section~\ref{sec:setting} fixes the algebraic setting and the
correspondence between constacyclic codes and labeled factor selections,
Section~\ref{sec:duality} develops component duality and the reciprocal-orbit
description of hull support, Section~\ref{sec:enumeration} derives the
enumerators with their statistical consequences, Section~\ref{sec:examples}
presents the examples and the computational validation, and
Section~\ref{sec:conclusion} collects the conclusions.

\begin{figure}[htbp]
\centering
\begin{tikzpicture}[
  node distance=2.4mm,
  stage/.style={draw,line width=0.4pt,align=center,inner sep=1.0mm,
                text width=20mm,minimum height=14mm,font=\scriptsize},
  outcome/.style={stage,line width=0.9pt},
  ar/.style={-{Latex[length=1.1mm,width=0.9mm]},line width=0.4pt}]
\node[stage]   (s1) {square-free\\affine\\algebra\\$\A\cong\prod_{s}K_s$};
\node[stage,right=of s1] (s2) {component\\moduli $M_s$\\$=x^n-\lambda_s$\\factors $\Fs$};
\node[stage,right=of s2] (s3) {compatible\\reciprocal\\$\tau_{s,k}$; orbits\\$O$, $a_O$, $w_O$};
\node[stage,right=of s3] (s4) {selection\\words $\varepsilon_O$\\code dim.\ $\kappa$,\\hull dim.\ $\eta_k$};
\node[stage,right=of s4] (s5) {one-orbit\\polynomial\\$\operatorname{tr}(T_{w}^{a})$\\global $\Epoly$};
\node[outcome,right=of s5] (s6) {distributions\\LCD count,\\moments,\\verification};
\draw[ar] (s1) -- (s2);
\draw[ar] (s2) -- (s3);
\draw[ar] (s3) -- (s4);
\draw[ar] (s4) -- (s5);
\draw[ar] (s5) -- (s6);
\end{tikzpicture}
\caption{The workflow of the paper. Each box names one construction of
Sections~\ref{sec:setting}--\ref{sec:enumeration}, in the order in which it is
applied: from the CRT decomposition of the affine algebra to the one-orbit
polynomial and the global joint enumerator
$\Epoly(u,z)=\prod_{s}\prod_{O\in\Os}\operatorname{tr}(T_{w_O}(u,z)^{a_O})$. The
final box lists the consequences of \S\ref{sec:consequences} and the validation
of \S\ref{sec:verification}.}
\label{fig:pipeline}
\end{figure}

\section{The affine algebra, constacyclic codes and factor selections}\label{sec:setting}

This section fixes the algebraic setting: the square-free affine algebra and its
CRT decomposition (\S\ref{sec:setup}), the correspondence between constacyclic
codes and labeled factor selections (\S\ref{sec:selections}), the global
$k$-Galois pairing (\S\ref{sec:pairing}), and the dimension and support
conventions used throughout (\S\ref{sec:dimensions}).

\subsection{Square-free affine algebras and their CRT decomposition}\label{sec:setup}

Throughout, $p$ is prime and
\[
q=p^e,\qquad e\geq 1 .
\]
For each component index $s$, let
\[
K_s=\F_{q^{m_s}}=\F_{p^{d_s}},\qquad d_s=e m_s,\qquad m_s\geq 1 .
\]
The component labels are fixed once and for all, and $\A$ always denotes the
ambient affine algebra.

The affine algebra is presented as
\begin{equation}\label{eq:algebra}
\A=\F_q[X_1,\ldots,X_\ell]
 \big/\langle t_1(X_1),\ldots,t_\ell(X_\ell)\rangle,
\end{equation}
where $\ell\geq 1$ and each $t_i$ is monic, square-free, and of positive degree.
Since each $t_i$ is square-free, $\A$ is reduced.

\begin{proposition}[Square-free component decomposition]\label{prop:decomposition}
The affine algebra \eqref{eq:algebra} has a fixed labeled decomposition
\begin{equation}\label{eq:decomp}
\A\cong\prod_{s=1}^{N}K_s,
\end{equation}
with finite fields $K_s=\F_{q^{m_s}}$.
\end{proposition}

\begin{proof}
Each $t_i$ is separable, so the one-variable quotient $\F_q[X_i]/\langle t_i\rangle$
is a finite \'etale $\F_q$-algebra: factoring $t_i$ into pairwise coprime monic
irreducibles and applying the Chinese remainder theorem exhibits it as a finite
product of finite separable field extensions of $\F_q$. The algebra
\eqref{eq:algebra} is the tensor product over $\F_q$ of these one-variable
quotients, and a tensor product of finite \'etale algebras is finite \'etale. A
finite \'etale algebra over a finite field is a finite product of finite fields.
Fixing the ordered tensor factors and the resulting CRT idempotents fixes the
component labels. The positive-degree hypothesis excludes the zero quotient.
\end{proof}

We fix the labeled isomorphism \eqref{eq:decomp} of
Proposition~\ref{prop:decomposition} throughout. Thus an element of $\A$ is
written as a tuple $(a_s)_s$, the primitive product idempotents are fixed by the
component labels, and all global constructions below use these tuple
coordinates. Dimensions of global objects are taken over $\F_q$, whereas
component dimensions are taken over $K_s$.

\subsection{Constacyclic codes and labeled factor selections}\label{sec:selections}

Fix $n\geq 1$ with $\gcd(n,p)=1$ and choose $\lambda_s\in K_s^*$ for every
component. Set
\[
M_s(x)=x^n-\lambda_s .
\]
Since $\gcd(n,p)=1$, each $M_s$ is square-free. Let $\Fs$ be the set of distinct
monic irreducible factors of $M_s$ over $K_s$. For a subset $J_s\subseteq\Fs$,
define
\[
g_{J_s}(x)=\prod_{f\in J_s}f(x),
\qquad
h_{J_s}(x)=\frac{M_s(x)}{g_{J_s}(x)} .
\]
The component code $C_s(J_s)$ is the principal ideal generated by $g_{J_s}$ in
$K_s[x]/\langle M_s\rangle$. Throughout, a factor is \emph{selected} if it is a
factor of the generator $g_{J_s}$, and we write $1$ for selected and $0$ for not
selected; a selected factor contributes to the codimension, not to the code
dimension.

\begin{lemma}[Factor-selection parametrization]\label{lem:factor-selection}
For each $s$, the distinct constacyclic ideals in the simple-root quotient
$K_s[x]/\langle M_s\rangle$ are in bijection with subsets $J_s\subseteq\Fs$.
Their $K_s$-dimension is
\[
\dim_{K_s}C_s(J_s)=n-\sum_{f\in J_s}\deg(f).
\]
The global product code has $\F_q$-dimension
\[
\dim_{\F_q}C(J)=\sum_{s=1}^{N}m_s\dim_{K_s}C_s(J_s).
\]
\end{lemma}

\begin{proof}
The simple-root quotient is a product of fields indexed by the distinct
irreducible factors of $M_s$. A principal ideal is obtained by choosing which
factor components vanish, equivalently by choosing the product $g_{J_s}$. Unique
factorization gives injectivity of the selection-to-ideal map, and every ideal
is obtained in this way. The usual degree formula gives the component dimension.
Since $[K_s:\F_q]=m_s$, the global dimension is the weighted sum displayed
above.
\end{proof}

The global selection space is therefore
\[
\mathscr C=\prod_{s=1}^{N}\{J_s:J_s\subseteq\Fs\},
\qquad
|\mathscr C|=2^{\sum_s|\Fs|} .
\]
This is a labeled factor-selection count: no quotient by rotations,
automorphisms, isometries, or code equivalence is taken, and the count does not
depend on the order of the component labels.

The component twists assemble into the global tuple
\[
\lambda=(\lambda_s)_{s=1}^{N}\in\A^* .
\]
The global constacyclic ambient module is
\begin{equation}\label{eq:ambient}
\mathcal M_\lambda
 =\A[x]/\langle x^n-\lambda\rangle
 \cong\prod_{s=1}^{N}K_s[x]/\langle x^n-\lambda_s\rangle
 \cong\prod_{s=1}^{N}K_s^n,
\end{equation}
where the last identification is the coefficient-vector identification. The
global constacyclic shift is therefore componentwise, and the fixed CRT
isomorphism embeds a global code in $\A^n$ as a tuple of component codes.

\begin{definition}[Constacyclic code over $\A$]\label{def:constacyclic-A}
Let $\lambda\in\A^{*}$. The \emph{$\lambda$-shift} on $\A^{n}$ is
\[
T_\lambda(c_0,c_1,\dots,c_{n-1})=(\lambda c_{n-1},c_0,\dots,c_{n-2}).
\]
A \emph{$\lambda$-constacyclic code of length $n$ over $\A$} is an
$\A$-submodule $C\subseteq\A^{n}$ with $T_\lambda(C)\subseteq C$. Under the
coefficient-vector identification $\A^{n}\cong\mathcal M_\lambda=\A[x]/\langle x^{n}-\lambda\rangle$,
$(c_0,\dots,c_{n-1})\mapsto\sum_i c_ix^{i}$, the map $T_\lambda$ becomes
multiplication by $x$. Hence the $\lambda$-constacyclic codes over $\A$ are
exactly the $\A[x]$-submodules (equivalently, the ideals) of
$\mathcal M_\lambda$.
\end{definition}

\begin{remark}\label{rem:A-linear}
The definition requires $\A$-linearity. A subspace of $\A^{n}$ that is only
$\F_q$-linear and closed under $T_\lambda$ need not be an $\A$-submodule, and
such subspaces are not counted here.
\end{remark}

For $J=(J_s)_s\in\mathscr C$ put
$C(J)=\prod_{s=1}^{N}C_s(J_s)\subseteq\prod_{s=1}^{N}K_s^n\cong\A^n$;
Lemma~\ref{lem:global-code} below shows that $C(J)$ is a $\lambda$-constacyclic
code over $\A$.

\begin{lemma}[Completeness of factor selections]\label{lem:complete}
Let $e_1,\dots,e_N$ be the primitive idempotents of $\A$ corresponding to the
fixed decomposition \eqref{eq:decomp} (so $e_s$ is the tuple with $1$ in
component $s$ and $0$ elsewhere). Every $\lambda$-constacyclic code
$C\subseteq\A^{n}$, that is, every $\A[x]$-submodule $C$ of
$\mathcal M_\lambda$, satisfies
\[
C=\bigoplus_{s=1}^{N} e_sC,
\]
where each $e_sC$ is an ideal of
$e_s\mathcal M_\lambda\cong K_s[x]/\langle x^{n}-\lambda_s\rangle$.
Consequently there is a unique $J=(J_s)_s\in\mathscr C$ with $C=C(J)$.
\end{lemma}

\begin{proof}
\emph{Step 1 (splitting).} The elements $e_s\in\A$ satisfy $e_se_t=\delta_{st}e_s$
and $\sum_s e_s=1$. Since $C$ is an $\A$-module, $e_sC\subseteq C$ for every $s$.
Any $c\in C$ can be written as $c=\sum_s e_sc$, so $C=\sum_s e_sC$. The sum is
direct: if $\sum_s e_sc_s=0$ with $c_s\in C$, multiplying by $e_t$ gives
$e_tc_t=0$ for every $t$.

\emph{Step 2 (components).} Multiplication by $e_s$ gives a surjective
$\A[x]$-linear map $\mathcal M_\lambda\to e_s\mathcal M_\lambda$. We have
$e_s\A=K_s$ and $e_s(x^{n}-\lambda)=e_s(x^{n}-\lambda_s)$. Therefore
\[
e_s\mathcal M_\lambda= e_s\A[x]\big/\langle e_s(x^{n}-\lambda)\rangle
 \cong K_s[x]/\langle x^{n}-\lambda_s\rangle,
\]
and this is the $s$-th factor of \eqref{eq:ambient}. The set $e_sC$ is the image
of an $\A[x]$-submodule under an $\A[x]$-linear map, so it is an
$\A[x]$-submodule of $e_s\mathcal M_\lambda$. On $e_s\mathcal M_\lambda$ the ring
$\A$ acts through the projection $\A\to K_s$. So $e_sC$ is a $K_s[x]$-submodule
of $K_s[x]/\langle M_s\rangle$, which is the same thing as an ideal.

\emph{Step 3 (identification).} By Lemma~\ref{lem:factor-selection}, every ideal
of the simple-root quotient $K_s[x]/\langle M_s\rangle$ equals
$C_s(J_s)=\langle g_{J_s}\rangle$ for exactly one $J_s\subseteq\Fs$. Applying
this to each $e_sC$ and using Step 1 gives $C=\bigoplus_s C_s(J_s)=C(J)$ with
$J=(J_s)_s$.

\emph{Uniqueness.} If $C(J)=C(J')$, multiplying by $e_s$ gives
$C_s(J_s)=C_s(J'_s)$ for every $s$, and Lemma~\ref{lem:factor-selection} gives
$J_s=J'_s$ for every $s$.
\end{proof}

\begin{lemma}[Global CRT ambient module and factor-selection codes]\label{lem:global-code}
Recall the definition of $C(J)$ for $J\in\mathscr C$. Equivalently, $C(J)$ is the
ideal of $\mathcal M_\lambda$ generated by the tuple $g_J=(g_{J_s})_s$. By
Definition~\ref{def:constacyclic-A}, each $C(J)$ is a $\lambda$-constacyclic code
over $\A$, and different labeled tuples $J\neq J'$ give different codes.
Together with Lemma~\ref{lem:complete}, the map
\[
\mathscr C\longrightarrow\{\lambda\text{-constacyclic codes of length }n\text{ over }\A\},
\qquad J\longmapsto C(J),
\]
is a bijection.
\end{lemma}

\begin{proof}
The product and coefficient-vector identifications are the CRT isomorphisms
\eqref{eq:ambient}. Each $C_s(J_s)$ is an ideal of $K_s[x]/\langle M_s\rangle$, so
$C(J)$ is an ideal of $\mathcal M_\lambda$. It is the ideal generated by $g_J$, and
it is a $\lambda$-constacyclic code over $\A$ by
Definition~\ref{def:constacyclic-A}. If $J\neq J'$, choose $s$ with
$J_s\neq J'_s$. Lemma~\ref{lem:factor-selection} gives
$C_s(J_s)\neq C_s(J'_s)$, so the two products have different $s$-components
(injectivity). Surjectivity is Lemma~\ref{lem:complete}.
\end{proof}

\subsection{The \texorpdfstring{$k$}{k}-Galois pairing and the global code}\label{sec:pairing}

For $x=(x_i)$ and $y=(y_i)$ in $K_s^n$, define
\[
\sigma_{s,k}(a)=a^{p^k},
\qquad
\rho_{s,k}(a)=a^{p^{d_s-k}}=\sigma_{s,k}^{-1}(a),
\qquad 0\leq k<e .
\]
When the component index is clear we write $\sigma(a)=a^{p^k}$ and
$\rho(a)=a^{p^{d_s-k}}$. The $k$-Galois pairing is
\begin{equation}\label{eq:pairing}
\langle x,y\rangle_{s,k}
 =\sum_{i=0}^{n-1}x_i\sigma_{s,k}(y_i)
 =\sum_{i=0}^{n-1}x_i y_i^{p^k},
\end{equation}
and we abbreviate $\langle x,y\rangle_k=\sum_i x_i y_i^{p^k}$ when the component
is fixed. The $k$-Galois dual and hull are
\[
\Dk(C)=\{y\in K_s^n:\langle c,y\rangle_{s,k}=0\text{ for every }c\in C\},
\qquad
\Hull_k(C)=C\cap \Dk(C).
\]
The pairing \eqref{eq:pairing} is $K_s$-linear in its first argument and
$\sigma_{s,k}$-semilinear in its second argument; the defining annihilator is
nevertheless a $K_s$-linear subspace. For $k=0$ it is Euclidean. On
$K_s=\F_{p^{d_s}}$ it is Hermitian precisely when $d_s$ is even and $k=d_s/2$;
thus for even $e$ the base-field choice $k=e/2$ is Hermitian on the components
with $m_s=1$ but not on those with $m_s>1$.

\begin{proposition}[$k$-Galois dual translation]\label{prop:dual-translation}
For every $K_s$-linear code $C\subseteq K_s^n$,
\[
\Dk(C)=\rho_{s,k}(C^{\perpE}),
\]
where $C^{\perpE}$ is the ordinary Euclidean dual and $\rho_{s,k}$ acts
coordinatewise.
\end{proposition}

\begin{proof}
The $k$-Galois condition is $\sum_i c_i\sigma_{s,k}(y_i)=0$ for every $c\in C$.
Setting $z=\sigma_{s,k}(y)$ gives $z\in C^{\perpE}$, and applying the inverse
coordinatewise map gives $y=\rho_{s,k}(z)$.
\end{proof}

\begin{lemma}[Global Frobenius and $k$-Galois pairing]\label{lem:global-pairing}
Define the componentwise maps on $\A$ by
\[
\sigma_{\A,k}((a_s)_s)=(\sigma_{s,k}(a_s))_s,
\qquad
\rho_{\A,k}((a_s)_s)=(\rho_{s,k}(a_s))_s .
\]
They are mutually inverse ring automorphisms. For $b\in\F_q$ and $a\in\A$, they
satisfy
\[
\sigma_{\A,k}(ba)=b^{p^k}\sigma_{\A,k}(a),\qquad
\rho_{\A,k}(ba)=b^{p^{e-k}}\rho_{\A,k}(a).
\]
Thus they are semilinear over $\F_q$ with respect to inverse base-field
automorphisms. On $\A^n$, define the $\A$-valued global $k$-Galois pairing
\[
\langle x,y\rangle_{\A,k}
 =\sum_{i=0}^{n-1}x_i\sigma_{\A,k}(y_i)
 =\left(\sum_{i=0}^{n-1}x_{s,i}y_{s,i}^{p^k}\right)_{s=1}^{N}.
\]
For a global code $C\subseteq\A^n$, define
\[
\begin{gathered}
D_{\A,k}(C)=\{y\in\A^n:\langle c,y\rangle_{\A,k}=0_{\A}\text{ for every }c\in C\},
\\
\Hull_{\A,k}(C)=C\cap D_{\A,k}(C).
\end{gathered}
\]
\end{lemma}

\begin{proof}
Each $\sigma_{s,k}$ is a field automorphism of $K_s$ with inverse $\rho_{s,k}$,
and the product of component automorphisms is a ring automorphism of $\A$. Under
the fixed CRT isomorphism, $\sigma_{\A,k}(a)=a^{p^k}$, so this is the intrinsic
power map on $\A$. The product idempotents are fixed because their coordinates are
$0$ or $1$. The restrictions of $\sigma_{s,k}$ and $\rho_{s,k}$ to $\F_q$ are
$b\mapsto b^{p^k}$ and $b\mapsto b^{p^{e-k}}$, respectively, since $d_s=em_s$.
The displayed pairing is the componentwise tuple of the already-defined component
pairings; zero in $\A\cong\prod_sK_s$ means zero in every component.
\end{proof}

\begin{lemma}[Global $k$-Galois dual decomposition]\label{lem:global-dual}
For every labeled factor-selection code,
\[
D_{\A,k}(C(J))
 =\prod_{s=1}^{N}\Dk(C_s(J_s)).
\]
\end{lemma}

\begin{proof}
Let $y=(y_s)_s$. If $y\in D_{\A,k}(C(J))$, fix $s$ and choose a global codeword
whose $s$-component is an arbitrary $c_s\in C_s(J_s)$ and whose other components
are zero. The $s$-component of the global pairing is then
$\langle c_s,y_s\rangle_{s,k}$, so it is zero for every $c_s$; hence
$y_s\in \Dk(C_s(J_s))$. This proves one inclusion. Conversely, if every $y_s$
belongs to its component dual, then every component of
$\langle c,y\rangle_{\A,k}$ is zero for every $c\in C(J)$, so the tuple pairing
is $0_{\A}$.
\end{proof}

\begin{lemma}[Global hull decomposition]\label{lem:global-hull}
The global $k$-Galois hull decomposes as
\[
\Hull_{\A,k}(C(J))
 =\prod_{s=1}^{N}\Hull_k(C_s(J_s)).
\]
\end{lemma}

\begin{proof}
Use the preceding dual decomposition and the elementary identity
$(\prod_s U_s)\cap(\prod_s V_s)=\prod_s(U_s\cap V_s)$ for subspaces in a finite
direct product.
\end{proof}

\subsection{Global dimensions, factor supports and notation}\label{sec:dimensions}

\begin{proposition}[Global $\F_q$-dimension additivity]\label{prop:global-dimensions}
For every $J\in\mathscr C$,
\[
\begin{gathered}
\dim_{\F_q}C(J)=\sum_{s=1}^{N}m_s\dim_{K_s}C_s(J_s),
\\
\dim_{\F_q}\Hull_{\A,k}(C(J))
 =\sum_{s=1}^{N}m_s\dim_{K_s}\Hull_k(C_s(J_s)).
\end{gathered}
\]
In particular, an orbit factor of common $K_s$-degree $d_O$ contributes the
global weight $w_O=m_s d_O$.
\end{proposition}

\begin{proof}
The $\F_q$-dimension of a finite direct product is the sum of the component
$\F_q$-dimensions, and $[K_s:\F_q]=m_s$. Apply the component-to-base-field
conversion to the global code and to the global hull decomposition. A factor of
$K_s$-degree $d_O$ contributes $m_s d_O$ over $\F_q$.
\end{proof}

Nondegeneracy of the component pairings implies that a code and its dual have
complementary global $\F_q$-dimensions. Hence the hull dimension cannot exceed
either the code dimension or its codimension in $\A^n$.

For a labeled selection define its generator-factor support and its surviving
CRT support by
\[
\operatorname{Supp}_{\mathrm{gen}}(J)=\{(s,f):f\in J_s\},
\qquad
\operatorname{Supp}_{\mathrm{CRT}}(C(J))
 =\{(s,f):f\in\Fs\setminus J_s\}.
\]
For any component or global ideal, $\operatorname{Supp}_{\mathrm{CRT}}$ denotes
the set of factor components on which the ideal is nonzero. The global hull
support is the union of the component hull supports of
Lemma~\ref{lem:global-hull}, and its $\F_q$-dimension is the sum in
Proposition~\ref{prop:global-dimensions}.

The notation of the paper is collected in Table~\ref{tab:nomenclature}; symbols
are used only in the sense fixed there.

\begin{table}[t]
\centering
\caption{Notation. All dimensions of global objects are over $\F_q$ and all
dimensions of component objects over $K_s$.}\label{tab:nomenclature}
\small
\begin{tabularx}{\textwidth}{@{}lX@{}}
\toprule
$q=p^{e}$, $n$, $k$ & base-field order, code length with $\gcd(n,p)=1$, and
twist exponent with $0\le k<e$ \\
$K_s=\F_{q^{m_s}}$, $d_s=em_s$ & component field and its absolute degree \\
$m_s$, $\Fs$, $J_s$, $\mathscr C$ & extension degree, irreducible factors of
$M_s=x^n-\lambda_s$, a labeled selection, the selection space \\
$\sigma_{s,k}$, $\rho_{s,k}$ & $a\mapsto a^{p^{k}}$ and its inverse
$a\mapsto a^{p^{d_s-k}}$ \\
$\langle x,y\rangle_{s,k}$, $\Dk$ & $k$-Galois pairing and $k$-Galois dual \\
$\A$ & square-free affine algebra, a labeled product of the $K_s$ \\
$g_{J_s}$, $h_{J_s}$ & selected-factor generator of $C_s(J_s)$ and
$M_s/g_{J_s}$ \\
$\lambda_s^{1+p^{d_s-k}}=1$ & compatibility condition \eqref{eq:compat} \\
$\tau_{s,k}$, $\Os$ & reciprocal factor permutation; its set of orbits
(Definition~\ref{def:orbits}) \\
$a_O$, $d_O$, $w_O=m_sd_O$ & orbit length, common factor degree, $\F_q$-weight \\
$\varepsilon_{O,i}$, $b_O$ & selection indicator; cyclic $1$-to-$0$ count
\eqref{eq:boundary} \\
$T_w(u,z)$ & weighted transfer matrix (rows: source state, columns: destination
state) \\
$\kappa(C)$, $\eta_k(C)$ & $\F_q$-dimension of a code and of its $k$-Galois hull \\
$\Epoly(u,z)$, $\mathcal H(z)$, $\Phi_j$ & joint enumerator, hull polynomial,
and the coefficient sums of Corollary~\ref{cor:conditional} \\
\bottomrule
\end{tabularx}
\end{table}

\section{\texorpdfstring{$k$}{k}-Galois duality and hull structure}\label{sec:duality}

We derive the $k$-Galois duals and hull supports of the component constacyclic
codes. Three objects are constructed: the inverse-Frobenius reciprocal and its
compatibility condition (\S\ref{sec:reciprocal}), the dual generator
(\S\ref{sec:duals}), and the hull support as a cyclic boundary statistic on
reciprocal orbits (\S\ref{sec:support}).

\subsection{The inverse-Frobenius reciprocal and the compatibility condition}\label{sec:reciprocal}

Let $f(x)=\sum_{i=0}^{r}f_i x^i$ be monic with $f_0\neq 0$. Define the normalized
inverse-Frobenius reciprocal
\begin{equation}\label{eq:reciprocal}
f^{\#_{s,k}}(x)
 =\rho_{s,k}(f_0)^{-1}
   \sum_{i=0}^{r}\rho_{s,k}(f_i)x^{r-i}.
\end{equation}

\begin{lemma}[Normalized inverse-Frobenius reciprocal]\label{lem:reciprocal}
The map $f\mapsto f^{\#_{s,k}}$ is monic, preserves degree and irreducibility, is
multiplicative on monic polynomials with nonzero constant term, and is inverted
by the corresponding normalized reciprocal using $\sigma_{s,k}$.
\end{lemma}

\begin{proof}
The leading coefficient after reversal is $\rho_{s,k}(f_0)\ne0$, so the displayed
normalization makes the reciprocal monic and of degree $r$. Its constant
coefficient is $\rho_{s,k}(f_r)/\rho_{s,k}(f_0)=\rho_{s,k}(f_0)^{-1}\ne0$.
Reversal satisfies $\operatorname{rev}(fg)=\operatorname{rev}(f)\operatorname{rev}(g)$,
and the coefficient map $\rho_{s,k}$ is multiplicative; the two normalization
factors also multiply, proving multiplicativity on monic polynomials.

For the inverse, write $g=f^{\#_{s,k}}$ and apply the corresponding normalized
reciprocal using $\sigma_{s,k}$. Since $g_0=\rho_{s,k}(f_0)^{-1}$, the inverse
normalization is $\sigma_{s,k}(g_0)^{-1}=f_0$, and coefficient reversal returns
$f$. The coefficient automorphism preserves irreducible factorizations, while
ordinary reciprocal reversal sends a nontrivial factorization of a polynomial
with nonzero constant term to a nontrivial factorization of its reciprocal. Thus
irreducibility is preserved.
\end{proof}

\begin{lemma}[Root action]\label{lem:root-action}
Let $f(x)=\sum_{i=0}^{r}f_ix^{i}\in K_s[x]$ be monic with $f_0\neq0$, let
$L\supseteq K_s$ be a finite field containing a root $\alpha$ of $f$, and let
$g=f^{\#_{s,k}}$ be as in \eqref{eq:reciprocal}. Then
$g\bigl(\alpha^{-p^{d_s-k}}\bigr)=0$. Thus the induced root map is
$\alpha\mapsto\alpha^{-p^{d_s-k}}$.
\end{lemma}

\begin{proof}
Extend $\rho=\rho_{s,k}$ to $L$ by $\rho(a)=a^{p^{d_s-k}}$; this is a field
automorphism of $L$ whose restriction to $K_s$ is $\rho_{s,k}$. Let
$\sigma_L:=\rho^{-1}$ on $L$; its restriction to $K_s$ is
$\sigma_{s,k}=\rho_{s,k}^{-1}$. For $y\in L$ put
\[
(\rho f)(y):=\rho\bigl(f(\sigma_L(y))\bigr)
 =\sum_{i=0}^{r}\rho(f_i)\,\rho(\sigma_L(y))^{i}
 =\sum_{i=0}^{r}\rho(f_i)\,y^{i},
\]
using that $\rho$ is a ring homomorphism. By \eqref{eq:reciprocal},
\[
g(x)=\rho(f_0)^{-1}\sum_{i=0}^{r}\rho(f_i)x^{r-i}
 =\rho(f_0)^{-1}x^{r}\,(\rho f)(x^{-1}).
\]
Put $\beta=\rho(\alpha^{-1})=\alpha^{-p^{d_s-k}}$, which is nonzero since
$f_0\ne0$ forces $\alpha\ne0$. Then $\beta^{-1}=\rho(\alpha)$, so
$\sigma_L(\beta^{-1})=\alpha$, and
\[
g(\beta)=\rho(f_0)^{-1}\beta^{r}\,(\rho f)(\beta^{-1})
 =\rho(f_0)^{-1}\beta^{r}\,\rho\bigl(f(\alpha)\bigr)=0.
\]
Finally, if $\alpha^{n}=\lambda_s$ then
$\beta^{n}=\rho(\alpha^{n})^{-1}=\rho(\lambda_s)^{-1}=\lambda_s^{-p^{d_s-k}}$.
\end{proof}

If $\alpha^n=\lambda_s$, then
\[
\bigl(\alpha^{-p^{d_s-k}}\bigr)^n=\lambda_s^{-p^{d_s-k}} .
\]
Consequently, the reciprocal preserves the same factor set precisely under
\begin{equation}\label{eq:compat}
\lambda_s^{1+p^{d_s-k}}=1 .
\end{equation}
Under this condition, define
\[
\tau_{s,k}:\Fs\longrightarrow\Fs,
\qquad
\tau_{s,k}(f)=f^{\#_{s,k}} .
\]
The map $\tau_{s,k}$ is a permutation. It need not be an involution; fixed points
and cycles of arbitrary admissible length are allowed.

\begin{proposition}[Compatible factor permutation]\label{prop:compatible-permutation}
Under $\lambda_s^{1+p^{d_s-k}}=1$, the map $\tau_{s,k}$ is a permutation of $\Fs$
and preserves factor degrees. Its cycles therefore partition the labeled factor
set into orbits with a common degree on each orbit.
\end{proposition}

\begin{proof}
The root action of Lemma~\ref{lem:root-action} sends roots of $M_s$ to roots of
the same polynomial under the displayed compatibility condition. By
Lemma~\ref{lem:reciprocal}, the normalized reciprocal \eqref{eq:reciprocal}
preserves monicity, irreducibility, and degree, and its inverse is obtained with
$\sigma_{s,k}$; hence it permutes $\Fs$.
\end{proof}

\subsection{\texorpdfstring{$k$}{k}-Galois duals}\label{sec:duals}

We first record the classical Euclidean dual of a constacyclic code, which is the
starting point for the dual generator.

\begin{lemma}[Ordinary constacyclic reciprocal]\label{lem:ordinary-reciprocal}
Let $K$ be a finite field, $a\in K^*$, and $n\geq1$. Let $M(x)=x^n-a$, let
$G\mid M$ be monic, and put $H=M/G$. If
\[
H^*(x)=H(0)^{-1}\sum_{i=0}^{\deg H}H_i x^{\deg H-i},
\]
then the Euclidean dual of the ideal $\langle G\rangle$ in
$K[x]/\langle M\rangle$ is the ideal generated by $H^*$ in the
$a^{-1}$-constacyclic quotient.
\end{lemma}

\begin{proof}
Write $r=\deg G$ and $h=\deg H$, so $r+h=n$. For the non-wrapped generator shifts
$x^iG$ and $x^jH^*$, with $0\leq i<h$ and $0\leq j<r$, the coefficient pairing
is, up to the nonzero scalar $H(0)^{-1}$, the coefficient of degree $h-i+j$ in
$GH=x^n-a$. Indeed, reversal changes the index from $u$ in $H$ to $h-u$ in $H^*$,
so the aligned terms satisfy $u+v=h-i+j$. Since $1\leq h-i+j\leq n-1$, that
coefficient is zero. These shifts therefore span mutually orthogonal spaces.

The reciprocal identity
\[
H^*(x)=H(0)^{-1}x^hH(x^{-1})
\]
shows that $H^*$ divides the normalized reciprocal of $M=x^n-a$, namely
$x^n-a^{-1}$. Hence its ideal is closed under the $a^{-1}$-constacyclic boundary
shift; the non-wrapped shifts above form a basis of dimension $r$. This dimension
equals the codimension of $\langle G\rangle$, so the orthogonal space is exactly
the $a^{-1}$-constacyclic ideal generated by $H^*$.
\end{proof}

\begin{theorem}[$k$-Galois dual generator]\label{thm:dual-generator}
Let $C_s(J_s)=\langle g_{J_s}\rangle$ be a component code and let
$h_{J_s}=M_s/g_{J_s}$. Under the simple-root hypotheses, the $k$-Galois dual is
generated by the normalized inverse-Frobenius reciprocal
$h_{J_s}^{\#_{s,k}}$ and is constacyclic with twist
\[
\lambda_s^{-p^{d_s-k}} .
\]
If the compatibility condition holds, the set of irreducible factors in its
generator is
\[
\tau_{s,k}(\Fs\setminus J_s).
\]
\end{theorem}

\begin{proof}
Proposition~\ref{prop:dual-translation} gives
$\Dk(C_s)=\rho_{s,k}(C_s^{\perpE})$. By
Lemma~\ref{lem:ordinary-reciprocal}, $C_s^{\perpE}$ is the
$\lambda_s^{-1}$-constacyclic ideal generated by the ordinary normalized
reciprocal $h_{J_s}^*$. Coordinatewise application of $\rho_{s,k}$ maps this
ideal onto the ideal generated by
$\rho_{s,k}(h_{J_s}^*)=h_{J_s}^{\#_{s,k}}$ in the quotient with twist
$\rho_{s,k}(\lambda_s^{-1})=\lambda_s^{-p^{d_s-k}}$. This calculation remains
valid without compatibility. When compatibility holds, multiplicativity and
Lemma~\ref{lem:root-action} map the irreducible factors of $h_{J_s}$, namely
$\Fs\setminus J_s$, onto $\tau_{s,k}(\Fs\setminus J_s)$; these are the selected
factors in the dual generator.
\end{proof}

\begin{proposition}[Same-twist compatibility criterion]\label{prop:same-twist}
The $k$-Galois dual twist equals the original twist precisely when
\[
\lambda_s^{-p^{d_s-k}}=\lambda_s,
\qquad\text{equivalently}\qquad
\lambda_s^{1+p^{d_s-k}}=1.
\]
If this condition fails, the $k$-Galois dual belongs to the transformed modulus
\[
x^n-\lambda_s^{-p^{d_s-k}},
\]
so the same-factor-set orbit product is not asserted.
\end{proposition}

\begin{proof}
The first equality is the twist in Theorem~\ref{thm:dual-generator}. Multiplying
it by $\lambda_s^{p^{d_s-k}}$ gives the equivalent displayed power condition. If
it fails, the root calculation places the dual in the transformed constacyclic
modulus rather than in the original same-factor-set quotient.
\end{proof}

\begin{remark}[Incompatible twists]\label{rem:incompatible}
If $\lambda_s^{1+p^{d_s-k}}\neq1$, put $\mu_s=\lambda_s^{-p^{d_s-k}}\ne\lambda_s$.
By Theorem~\ref{thm:dual-generator}, the dual generator $h_{J_s}^{\#_{s,k}}$
divides $x^n-\mu_s$, whereas $g_{J_s}$ divides $x^n-\lambda_s$. The two moduli
are coprime because their difference is the nonzero constant $\mu_s-\lambda_s$.
Consequently the two generators are coprime, and their degrees sum to $n$. A
vector in the hull, identified with a polynomial of degree less than $n$, must
be divisible by both generators and hence by their product of degree $n$.
Therefore
\[
\Hull_k(C_s(J_s))=\{0\}.
\]
Every such component code is $k$-Galois LCD. The same-factor-set permutation
$\tau_{s,k}$ is unavailable in this case, so the orbit product of
Theorem~\ref{thm:central} remains restricted to compatible twists.
\end{remark}

\begin{remark}[Dual convention]\label{rem:candidate}
The pairing \eqref{eq:pairing} is code-first: the Frobenius is applied to the
second argument. Some of the literature instead forms the dual with the
candidate-first condition
\[
D_{\mathrm{cand}}(C)=\Bigl\{y:\ \sum_i y_i\,\sigma_{s,k}(c_i)=0\text{ for every }c\in C\Bigr\}.
\]
\textup{(a)} For any $K_s$-linear code, with no constacyclic hypothesis,
\[
\Dk(C)=\rho_{s,k}(C^{\perpE}),\qquad
D_{\mathrm{cand}}(C)=\sigma_{s,k}(C^{\perpE}),
\]
and therefore
\begin{equation}\label{eq:convention}
D_{\mathrm{cand}}(C)=\sigma_{s,k}^{2}\bigl(\Dk(C)\bigr).
\end{equation}
The two dual subspaces are not generally equal.
\emph{Proof.} The first relation is Proposition~\ref{prop:dual-translation}.
For the candidate-first condition, substitute $y=\sigma_{s,k}(z)$: then
$\sum_iy_i\sigma_{s,k}(c_i)=\sigma_{s,k}(\sum_iz_ic_i)$, which vanishes for every
$c\in C$ exactly when $z\in C^{\perpE}$, giving
$D_{\mathrm{cand}}(C)=\sigma_{s,k}(C^{\perpE})$. Applying $\sigma_{s,k}$ twice to
the $k$-Galois relation gives \eqref{eq:convention}, since
$\sigma_{s,k}^{2}(\rho_{s,k}(z))=\sigma_{s,k}(z)$. The two subspaces agree only
when the square of the Frobenius happens to fix the dual subspace.

\textup{(b)} \emph{(Same-code hull dimension only.)} Let $B$ be a row basis of a
finite-field-linear code $C$ and let $G=B\,\sigma(B)^{\mathsf T}$. Then both
same-code hull dimensions equal $\dim(C)-\rank(G)$. In particular, the
$k$-Galois and candidate-first pairings give equal hull dimensions for the same
code, and the same LCD decision. This does not identify the dual codes, hull
subspaces, factor supports, or generator supports.
\emph{Proof.} Write codewords as $vB$ with row vectors $v$. For the $k$-Galois
hull, $y=vB$ lies in $\Dk(C)$ exactly when
$uB\,\sigma(B)^{\mathsf T}\sigma(v)^{\mathsf T}=0$ for every $u$, that is,
$G\,\sigma(v)^{\mathsf T}=0$. Since $\sigma$ is a coordinatewise field
automorphism, $\sigma(v)$ ranges over all row vectors as $v$ does, so the
$k$-Galois hull has dimension $\dim(C)-\rank(G)$ by rank-nullity. For the
candidate-first hull, $y=vB$ lies in $D_{\mathrm{cand}}(C)$ exactly when
$vB\,\sigma(B)^{\mathsf T}\sigma(u)^{\mathsf T}=0$ for every $u$, that is,
$vG=0$; this left nullspace has the same dimension $\dim(C)-\rank(G)$. The LCD
assertion is the case $\rank(G)=\dim(C)$.
\end{remark}

\subsection{Hull support and the cyclic boundary statistic}\label{sec:support}

\begin{definition}[Orbit set]\label{def:orbits}
Assume $\lambda_s^{1+p^{d_s-k}}=1$. The \emph{orbit set} $\Os$ is the set of
cycles of the permutation $\tau_{s,k}$ of $\Fs$
(Proposition~\ref{prop:compatible-permutation}). For $O\in\Os$ write
$a_O=|O|$, let $d_O$ be the common $K_s$-degree of its factors, and list
$O=\{f_{O,1},\dots,f_{O,a_O}\}$ so that
$\tau_{s,k}(f_{O,i})=f_{O,i+1}$, indices modulo $a_O$. Thus
$\Fs=\bigsqcup_{O\in\Os}O$.
\end{definition}

\begin{theorem}[Component hull support]\label{thm:hull-support}
Let $C_s(J_s)$ be selected by $J_s\subseteq\Fs$. Under the compatible
same-factor-set hypotheses, the $k$-Galois hull has factor support
\begin{equation}\label{eq:support}
\begin{aligned}
\operatorname{Supp}_{\mathrm{CRT}}\bigl(\Hull_k(C_s(J_s))\bigr)
 &= (\Fs\setminus J_s)\cap\tau_{s,k}(J_s)\\
 &= \tau_{s,k}(J_s)\setminus J_s.
\end{aligned}
\end{equation}
The first line is the direct lcm/intersection form, and the second uses that
$\tau_{s,k}$ is a permutation of $\Fs$.
\end{theorem}

\begin{proof}
The code generator has support $J_s$, while the dual generator has support
$\tau_{s,k}(\Fs\setminus J_s)$. The intersection of the two principal ideals in
the square-free quotient is generated by the least common multiple, so its
surviving factors are the complement of the union of these two supports. This
gives
$(\Fs\setminus J_s)\cap(\Fs\setminus\tau_{s,k}(\Fs\setminus J_s))$.
Since $\tau_{s,k}$ is a permutation, the second complement is $\tau_{s,k}(J_s)$,
giving the two displayed forms.
\end{proof}

\begin{proposition}[Weighted cyclic boundary statistic]\label{prop:boundary}
For a labeled selection $J=(J_s)_s$ and its global code $C(J)$, let $O\in\Os$ be
indexed as in Definition~\ref{def:orbits}. All factors in $O$ have a common
degree $d_O$. Put
\[
w_O=m_s d_O,
\qquad
\varepsilon_{O,i}=\begin{cases}1,&f_{O,i}\in J_s,\\0,&f_{O,i}\notin J_s,\end{cases}
\]
following the selection convention of \S\ref{sec:selections}; when the orbit is
clear we write $\varepsilon_i=\varepsilon_{O,i}$. Then
\begin{equation}\label{eq:boundary}
b_O(\varepsilon)=\sum_{i=1}^{a_O}\varepsilon_i(1-\varepsilon_{i+1})
\end{equation}
and
\begin{equation}\label{eq:hull-dim}
\dim_{\F_q}\Hull_{\A,k}(C(J))
 =\sum_{s=1}^{N}\sum_{O\in\Os}w_O b_O(\varepsilon_O).
\end{equation}
\end{proposition}

\begin{proof}
The orbit positions have equal factor degree because the reciprocal permutation
of Proposition~\ref{prop:compatible-permutation} preserves degree. The support
formula \eqref{eq:support} identifies a hull factor at position $i+1$ exactly
when $f_i$ is selected and $f_{i+1}$ is not. Each such factor contributes
$m_s d_O=w_O$ to the global $\F_q$ dimension, and summing over the orbit
partition gives the formula.
\end{proof}

\begin{proposition}[Component lcm--boundary bridge]\label{prop:lcm-bridge}
Under the compatible simple-root assumptions, write $g_s=g_{J_s}$,
$h_s=M_s/g_s$, and take the normalized $k$-Galois reciprocal
$h_s^{\#_{s,k}}$. Then the global $\F_q$-dimension satisfies
\begin{equation}\label{eq:lcm-bridge}
\begin{aligned}
\dim_{\F_q}\Hull_{\A,k}(C(J))
 &=\sum_{s=1}^{N}m_s\bigl(n-\deg\operatorname{lcm}(g_s,h_s^{\#_{s,k}})\bigr)\\
 &=\sum_{s,O}w_O\sum_i\varepsilon_{O,i}(1-\varepsilon_{O,i+1}).
\end{aligned}
\end{equation}
The factor $m_s$ is necessary when $K_s\ne\F_q$; replacing the first sum by a
single unweighted $n-\deg\operatorname{lcm}$ is generally incorrect.
\end{proposition}

\begin{proof}
In the square-free component modulus, the selected-factor exponent in $g_s$ is
$\varepsilon_f$, and that in the dual generator at $f$ is
$1-\varepsilon_{\tau_{s,k}^{-1}(f)}$. Therefore the lcm contains $f$ with
exponent $\max\{\varepsilon_f,1-\varepsilon_{\tau^{-1}(f)}\}$. Subtracting its
degree from $n=\sum_{f\in\Fs}\deg f$ leaves the factor weight
$\deg f\,\varepsilon_{\tau^{-1}(f)}(1-\varepsilon_f)$. Multiply by $m_s$, sum
over all components, and index each orbit in the forward direction
$\tau(f_i)=f_{i+1}$. This gives the second expression. The first expression is
the usual degree formula for an intersection of principal ideals over each $K_s$,
followed by extension-degree conversion. The weight $m_s=[K_s:\F_q]$ converts
each component $K_s$-dimension to an $\F_q$-dimension; it is essential whenever
$K_s\ne\F_q$, for instance in the $\F_4\times\F_{16}$ product of
\S\ref{sec:examples}, where the second component's hull dimensions double.
\end{proof}

The subfamily of \cite{debnathAffine2026} that overlaps the setting above is
$m_s=1$, $0\le k<e$, $\gcd(\operatorname{ord}(\lambda_s),n)=1$ and
$\operatorname{ord}(\lambda_s)\mid(1+p^{e-k})$. There, the hull dimension is
expressed by factor multiplicities $u_i$ and terms
$\max\{u_i,p^{e'}-u_{i-1}\}$, under the additional hypotheses $n=n'p^{e'}$,
$\gcd(n',p)=1$, $\gcd(\operatorname{ord}(\lambda),n')=1$ and
$\operatorname{ord}(\lambda)\mid(1+p^{e-k})$; the formula is Thm.~8 of the
preprint version (\href{https://arxiv.org/abs/2412.08512v1}{arXiv:2412.08512v1})
of \cite{debnathAffine2026}. At $e'=0$ the multiplicities are binary,
$u_i=\varepsilon_i\in\{0,1\}$, and the contribution complementary to each maximum
is
\[
1-\max\{\varepsilon_i,1-\varepsilon_{i-1}\}
 =\varepsilon_{i-1}(1-\varepsilon_i).
\]
After matching the factor order and degrees, this is precisely the corresponding
summand in \eqref{eq:lcm-bridge}, so the two descriptions agree on this
binary-max specialization on matched factor orderings. The present theorem
imposes no coprimality between $n$ and the order of the twist. The comparison
does not identify the candidate-first and $k$-Galois dual subspaces; their
same-code hull dimensions can be compared through Remark~\ref{rem:candidate}(b).

\begin{proposition}[Cyclic boundary interpretation]\label{prop:boundary-interpretation}
For a nonconstant cyclic binary word, the statistic $b_O$ of \eqref{eq:boundary}
equals the number of one-runs and also the number of zero-runs. It is the number
of directed $1$-to-$0$ transitions, with cyclic indexing modulo $a_O$.
\end{proposition}

\begin{proof}
Each summand is one exactly when the directed edge from $i$ to $i+1$ is a
$1$-to-$0$ transition. Every one-run has exactly one such exit and every zero-run
has exactly one entry, so these counts agree. The two constant words have
$b_O=0$.
\end{proof}

Figure~\ref{fig:orbit} illustrates one orbit with its selection word: the two
thick $1$-to-$0$ arrows are exactly the transitions counted by
\eqref{eq:boundary}.

\begin{figure}[ht]
\centering
\begin{tikzpicture}[>=stealth, box/.style={draw, minimum width=1.2cm,
  minimum height=1.0cm, align=center, inner sep=1.2pt}]
\node[box] (f0) at (0,0) {$f_{1}$\\$\varepsilon_{1}=1$};
\node[box] (f1) at (1.82,0) {$f_{2}$\\$\varepsilon_{2}=1$};
\node[box] (f2) at (3.64,0) {$f_{3}$\\$\varepsilon_{3}=0$};
\node[box] (f3) at (5.46,0) {$f_{4}$\\$\varepsilon_{4}=1$};
\node[box] (f4) at (7.28,0) {$f_{5}$\\$\varepsilon_{5}=0$};
\node[box] (f5) at (9.10,0) {$f_{6}$\\$\varepsilon_{6}=0$};
\draw[->] (f0) -- node[above] {\footnotesize$\tau$} (f1);
\draw[->, ultra thick] (f1) -- (f2);
\draw[->] (f2) -- (f3);
\draw[->, ultra thick] (f3) -- (f4);
\draw[->] (f4) -- (f5);
\draw[->] (f5.south) -- ++(0,-0.5) -| node[below, near start]
  {\footnotesize$\tau$} (f0.south);
\node[font=\small, align=center] at (4.55,-1.8) {thick arrows: the
  $1$-to-$0$ transitions, so $b_O=2$ and the hull contribution is $2w_O$};
\end{tikzpicture}
\caption{One reciprocal orbit $O$ of length $a_O=6$, with factors
$f_{O,1},\dots,f_{O,6}$ written $f_1,\dots,f_6$ and selection word
$\varepsilon=110100$. The arrows follow $\tau=\tau_{s,k}$; the two thick arrows
are the $1$-to-$0$ transitions counted by \eqref{eq:boundary}, and each of them
contributes one factor of $K_s$-degree $d_O$, hence $w_O=m_sd_O$ to the global
$\F_q$-hull dimension by \eqref{eq:hull-dim}.}
\label{fig:orbit}
\end{figure}

\section{Exact joint enumeration}\label{sec:enumeration}

By Theorem~\ref{thm:hull-support} and Proposition~\ref{prop:boundary}, both
statistics of interest are additive over reciprocal orbits. The code dimension
counts the unselected factors, weighted by $w_O$; the $k$-Galois hull dimension
counts the $1$-to-$0$ transitions of the selection word, again weighted by $w_O$.
The labeled selections therefore factorize over orbits, and the joint enumerator
is a product of independent one-orbit polynomials. This section
computes those polynomials in two equivalent ways
(\S\S\ref{sec:orbit}--\ref{sec:transfer}), assembles the global enumerator
(\S\ref{sec:global}), records its consequences (\S\ref{sec:consequences}) and
describes the exact coefficient computation (\S\ref{sec:algorithm}).

\subsection{The one-orbit polynomial}\label{sec:orbit}

\begin{proposition}[One-orbit polynomial]\label{prop:orbit-poly}
For an orbit of length $a\geq1$ and positive integer weight $w$, let
$P_{a,w}(z)$ be the sum of $z^{wb(\varepsilon)}$ over all indexed cyclic binary
words. For $1\leq r\leq\lfloor a/2\rfloor$,
\[
N(a,r)=2\binom{a}{2r}
 =\frac{a}{r}\binom{a-1}{2r-1}
\]
counts the nonconstant words with $r$ one-runs, and
\begin{equation}\label{eq:orbit-poly}
P_{a,w}(z)=2+
 \sum_{r=1}^{\lfloor a/2\rfloor}
 2\binom{a}{2r}z^{rw}
 =2+
 \sum_{r=1}^{\lfloor a/2\rfloor}
 \frac{a}{r}\binom{a-1}{2r-1}z^{rw}.
\end{equation}
\end{proposition}

\begin{proof}
The two constant words contribute $2$. By
Proposition~\ref{prop:boundary-interpretation}, a nonconstant word with $r$
one-runs has $2r$ transition edges. Choosing the $2r$ transition edges and one of
the two starting values gives $2\binom{a}{2r}$. The equality with the composition
form follows from $2\binom{a}{2r}=(a/r)\binom{a-1}{2r-1}$, so the second
expression is integral because it equals the first.
\end{proof}

For example,
\[
P_{1,w}(z)=2,\quad
P_{2,w}(z)=2+2z^w,\quad
P_{3,w}(z)=2+6z^w,
\]
\[
P_{4,w}(z)=2+12z^w+2z^{2w},\qquad
P_{5,w}(z)=2+20z^w+10z^{2w}.
\]
At $a=4$, for instance, there are two constant words, twelve words with one
one-run, and two alternating words with two one-runs. At $z=1$ the binomial
identity gives $P_{a,w}(1)=2^a$, as required for the $2^a$ indexed factor
selections on one orbit.

\subsection{Transfer matrices and the closed form}\label{sec:transfer}

The one-orbit polynomial admits a two-state transfer-matrix representation, whose
trace enumerates the closed walks that encode the cyclic words, and the closed
form \eqref{eq:orbit-poly}, which exhibits the coefficients directly. The two are
used for different purposes: the transfer matrix gives the uniform recursion of
\S\ref{sec:algorithm}, whereas \eqref{eq:orbit-poly} gives each coefficient in
closed binomial form.

\begin{proposition}[Weighted transfer matrix]\label{prop:transfer-matrix}
For a transition from source state $r=\varepsilon_i$ to destination state
$t=\varepsilon_{i+1}$, assign
\[
T_w(u,z)_{r,t}=u^{w(1-t)}z^{wr(1-t)}.
\]
With states ordered as $0,1$, this is
\begin{equation}\label{eq:transfer}
T_w(u,z)=\begin{pmatrix}u^w&1\\u^wz^w&1\end{pmatrix}.
\end{equation}
The destination weight is deliberate: each factor is counted once as the
destination of the preceding edge. Thus $u$ records code dimension, because a
factor with destination state $t=0$ is not selected and contributes $w$ to the
code dimension. The factor $z^w$ appears precisely on a $1$-to-$0$ edge.
\end{proposition}

\begin{proof}
The four source/destination choices $(r,t)\in\{0,1\}^2$ receive the destination
code-dimension weight $u^{w(1-t)}$ and the boundary weight $z^{wr(1-t)}$. This
gives the displayed matrix in the state order $0,1$.
\end{proof}

For illustration, the bivariate short-orbit traces are
\[
\operatorname{tr}(T_w(u,z))=1+u^w,\qquad
\operatorname{tr}(T_w(u,z)^2)=1+u^{2w}+2u^wz^w.
\]
A one-factor orbit never contributes a hull boundary; for length two, the two
mixed selections each contribute one boundary, not two.

\begin{proposition}[Local trace identity]\label{prop:trace}
With $T_w(u,z)$ as in Proposition~\ref{prop:transfer-matrix}, for an orbit of
length $a$ and weight $w$,
\begin{equation}\label{eq:trace}
\operatorname{tr}(T_w(u,z)^a)
 =\sum_{\varepsilon\in\{0,1\}^a}
 u^{w\sum_i(1-\varepsilon_i)}z^{w\sum_i\varepsilon_i(1-\varepsilon_{i+1})}.
\end{equation}
In particular,
\[
\operatorname{tr}(T_w(1,z)^a)=P_{a,w}(z).
\]
\end{proposition}

\begin{proof}
Expand a diagonal entry of $T_w^a$ in \eqref{eq:transfer} as a sum over
intermediate states. The diagonal condition identifies the terminal state with
the initial state, thereby closing the cycle. Summing the two diagonal entries
sums over every indexed binary word exactly once. The product of edge weights is
\[
\prod_i u^{w(1-\varepsilon_{i+1})}z^{w\varepsilon_i(1-\varepsilon_{i+1})},
\]
which is the displayed monomial after cyclic reindexing.
\end{proof}

\begin{proposition}[Closed form of the one-orbit trace]\label{prop:closed}
For an orbit of length $a\ge1$ and positive integer weight $w$, and for
$1\le j\le a-1$ and $1\le r\le\min\{j,a-j\}$, the number of indexed cyclic binary
words $\varepsilon\in\{0,1\}^{a}$ with exactly $j$ zeros (unselected factors) and
exactly $r$ one-runs is
\[
N(a;j,r)=\frac{a}{r}\binom{j-1}{r-1}\binom{a-j-1}{r-1}.
\]
Consequently
\[
\operatorname{tr}\bigl(T_w(u,z)^{a}\bigr)
 =1+u^{wa}+\sum_{j=1}^{a-1}\sum_{r=1}^{\min\{j,a-j\}}
 \frac{a}{r}\binom{j-1}{r-1}\binom{a-j-1}{r-1}\,u^{wj}z^{wr}.
\]
\end{proposition}

\begin{proof}
\emph{Double counting.} Fix $1\le j\le a-1$, so $\varepsilon$ is not constant,
and fix $1\le r\le\min\{j,a-j\}$. Each of the $r$ one-runs and $r$ zero-runs has
positive length, so no such words exist for larger $r$. Call $i\in\mathbb{Z}/a$ a
\emph{run start} of $\varepsilon$ if $\varepsilon_i=1$ and $\varepsilon_{i-1}=0$.
Each one-run of a nonconstant cyclic word has exactly one run start, so a word
with $r$ one-runs has exactly $r$ run starts. Let $\mathcal{P}$ be the set of
pairs $(\varepsilon,i)$ where $\varepsilon$ has $j$ zeros and $r$ one-runs and $i$
is a run start of $\varepsilon$. Then $|\mathcal{P}|=r\,N(a;j,r)$.

\emph{Bijection.} Send $(\varepsilon,i)$ to $(i,v)$ with
$v=(\varepsilon_i,\varepsilon_{i+1},\dots,\varepsilon_{i+a-1})$, the word read
linearly from position $i$. Then $v$ starts with $1$ (as $\varepsilon_i=1$) and
ends with $0$ (as $\varepsilon_{i-1}=0$). Hence $v$ is a linear word consisting of
$2r$ maximal runs alternating $1,0,1,0,\dots,1,0$: $r$ one-runs with total length
$a-j$ and $r$ zero-runs with total length $j$. (The cyclic runs of $\varepsilon$
are exactly the linear runs of $v$, because the cut between positions $i-1$ and
$i$ lies at a $0\to1$ transition.) Conversely, any pair $(i,v)$ with
$i\in\mathbb{Z}/a$ and $v$ of this shape gives, by placing $v$ cyclically starting
at position $i$, a word with $j$ zeros, $r$ one-runs and run start $i$. The two
maps are mutually inverse. Such words $v$ correspond to an ordered composition of
$a-j$ into $r$ positive parts and an ordered composition of $j$ into $r$ positive
parts, so their number is $\binom{a-j-1}{r-1}\binom{j-1}{r-1}$. Therefore
\[
r\,N(a;j,r)=|\mathcal{P}|=a\binom{j-1}{r-1}\binom{a-j-1}{r-1},
\]
which gives the formula.

\emph{Trace.} By Proposition~\ref{prop:trace}, $\operatorname{tr}(T_w^{a})$ is
the sum over $\varepsilon\in\{0,1\}^{a}$ of
$u^{w\cdot\#\{i:\varepsilon_i=0\}}\allowbreak z^{w\,b(\varepsilon)}$. By
Proposition~\ref{prop:boundary-interpretation}, $b(\varepsilon)$ is the number of
one-runs of a nonconstant word. The all-zero word gives $u^{wa}z^{0}$ and the
all-one word gives $u^{0}z^{0}=1$. Grouping the nonconstant words by $(j,r)$ gives
the stated sum.
\end{proof}

\begin{remark}
Summing over $j$ at fixed $r$ recovers $N(a,r)=2\binom{a}{2r}$, because both count
the words with $r$ one-runs (Proposition~\ref{prop:orbit-poly}).
\end{remark}

\subsection{The global joint enumerator}\label{sec:global}

It remains to multiply the orbit factors. The variable $u$ carries the exponent
$\kappa$ of the global code dimension and $z$ the exponent $\eta_k$ of the
$k$-Galois hull dimension. A factor of $K_s$-degree $d_O$ on an orbit of the
component $K_s$ contributes the weight $w_O=m_sd_O$ to both exponents. A labeled
selection is an independent choice of a binary word on each reciprocal orbit,
and both statistics are sums of the corresponding orbit contributions; hence the
bivariate generating polynomial of the population $(\kappa,\eta_k)$ is the
product of the one-orbit traces.

\begin{keyresult}
\begin{theorem}[Exact labeled joint $(\kappa,\eta)$-enumerator]\label{thm:central}
Under the standing hypotheses of Sections~\ref{sec:setting} and~\ref{sec:duality},
with the orbit notation of Definition~\ref{def:orbits} and
\S\S\ref{sec:orbit}--\ref{sec:transfer}, let
\[
\A\cong\prod_{s=1}^{N}K_s,\qquad K_s=\F_{q^{m_s}},\qquad
\lambda=(\lambda_s)_s\in\A^{*},
\]
with $q=p^{e}$, $d_s=em_s$, $0\le k<e$, $\gcd(n,p)=1$ (so every $M_s$ has simple
roots), fixed component labels, $\lambda$-constacyclic codes over $\A$ in the
sense of Definition~\ref{def:constacyclic-A}, and
\[
\lambda_s^{1+p^{d_s-k}}=1\qquad(1\le s\le N).
\]
For a labeled factor selection $J=(J_s)_s$, the global code is the CRT product
\[
C(J)=\prod_{s=1}^{N}C_s(J_s)\subseteq\A^{n},
\]
and its global $k$-Galois dual and hull are defined by the $\A$-valued pairing of
Lemma~\ref{lem:global-pairing}:
\[
\begin{gathered}
D_{\A,k}(C(J))=\{y:\langle c,y\rangle_{\A,k}=0_{\A}\ \text{for every } c\in C(J)\},
\\
\Hull_{\A,k}(C(J))=C(J)\cap D_{\A,k}(C(J)).
\end{gathered}
\]
For every orbit $O\in\Os$, let $a_O$ be its length, let $d_O$ be its common factor
degree, and put $w_O=m_sd_O$. With $\varepsilon_{O,i}$ as in
Proposition~\ref{prop:boundary}, set
\begin{equation}\label{eq:KH}
\begin{aligned}
\kappa(C(J)) &:= \dim_{\F_q}C(J)
 =\sum_{s,O,i}w_O(1-\varepsilon_{O,i}),\\
\eta_k(C(J)) &:= \dim_{\F_q}\Hull_{\A,k}(C(J))
 =\sum_{s,O,i}w_O\,\varepsilon_{O,i}(1-\varepsilon_{O,i+1}).
\end{aligned}
\end{equation}
With
\[
T_w(u,z)=\begin{pmatrix}u^{w}&1\\ u^{w}z^{w}&1\end{pmatrix},
\]
the exact joint generating polynomial is
\begin{equation}\label{eq:joint}
\begin{aligned}
\Epoly(u,z)
 &=\sum_{J\in\mathscr C}u^{\kappa(C(J))}z^{\eta_k(C(J))}\\
 &=\prod_{s=1}^{N}\prod_{O\in\Os}
 \operatorname{tr}\bigl(T_{w_O}(u,z)^{a_O}\bigr).
\end{aligned}
\end{equation}
For every pair $(\kappa,\eta)$ of \eqref{eq:KH}, $[u^{\kappa}z^{\eta}]\,\Epoly(u,z)$
is the number of $\lambda$-constacyclic codes $C\subseteq\A^{n}$ (equivalently, by
Lemmas~\ref{lem:complete} and~\ref{lem:global-code}, of labeled factor selections
$J\in\mathscr C$) with global $\F_q$-code dimension $\kappa$ and global
$k$-Galois $\F_q$-hull dimension $\eta$.
\end{theorem}
\end{keyresult}

\begin{proof}
\emph{Product structure.} By Lemma~\ref{lem:factor-selection}, a labeled
selection is a tuple $J=(J_s)_s$ with $J_s\subseteq\Fs$ arbitrary, so
$\mathscr C=\prod_s 2^{\Fs}$. By Proposition~\ref{prop:compatible-permutation},
under $\lambda_s^{1+p^{d_s-k}}=1$ the permutation $\tau_{s,k}$ partitions each
$\Fs$ into its orbits $O\in\Os$. With each orbit listed as in
Definition~\ref{def:orbits}, a subset $J_s\subseteq\Fs$ is the same as a family of
binary words $(\varepsilon_{O,i})_{i=1}^{a_O}$, one for each $O\in\Os$. Hence
\[
\mathscr C\;\cong\;\prod_{s=1}^{N}\prod_{O\in\Os}\{0,1\}^{a_O},
\]
and the word on each reciprocal-factor orbit is an independent coordinate.

\emph{Additivity.} Lemmas~\ref{lem:global-dual} and~\ref{lem:global-hull}
identify the global dual and hull with the products of the component duals and
hulls. Proposition~\ref{prop:global-dimensions} converts component dimensions to
global $\F_q$-dimensions with the weights $w_O=m_sd_O$. Hence both statistics in
\eqref{eq:KH} are sums over orbits: $\kappa$ by
Lemma~\ref{lem:factor-selection}, and $\eta_k$ by
Theorem~\ref{thm:hull-support} with Proposition~\ref{prop:boundary}.

\emph{Counting.} Proposition~\ref{prop:trace} enumerates each orbit coordinate
with its code-dimension and hull-dimension weights. A product of generating
polynomials convolves independent additive statistics, so the coefficient of
$u^{\kappa}z^{\eta}$ in the orbit-trace product counts exactly the tuples of orbit
words with total weights $(\kappa,\eta)$.

\emph{From selections to codes.} By Lemma~\ref{lem:global-code} (injectivity),
distinct tuples give distinct codes. By Lemma~\ref{lem:complete} (surjectivity),
every $\lambda$-constacyclic code over $\A$ arises from some tuple. So the
coefficient counts $\lambda$-constacyclic codes over $\A$ with parameters
$(\kappa,\eta)$.
\end{proof}

\subsection{Consequences}\label{sec:consequences}

Specialising \eqref{eq:joint} yields the marginal distributions, the LCD count and
the moments of the hull dimension, and differentiating it at fixed code dimension
yields the conditional hull statistics.

\begin{corollary}\label{cor:total}
The total number of $\lambda$-constacyclic codes of length $n$ over $\A$ is
\[
\Epoly(1,1)=2^{\sum_s|\Fs|}.
\]
\end{corollary}

\begin{proof}
At $u=z=1$, each local trace \eqref{eq:trace} counts the $2^{a_O}$ indexed binary
words on its orbit. Multiplying over the orbit partition gives
$2^{\sum_s|\Fs|}$. The same number follows directly from the bijection
$J\mapsto C(J)$ of Lemmas~\ref{lem:complete} and~\ref{lem:global-code} and
$|\mathscr C|=2^{\sum_s|\Fs|}$. (This count uses only
Lemma~\ref{lem:factor-selection} and does not need the compatibility condition.)
\end{proof}

\begin{corollary}[Code-dimension distribution]\label{cor:code-dist}
Setting $z=1$ removes the hull marking and yields
\begin{equation}\label{eq:dimdist}
\Epoly(u,1)=\prod_{s=1}^{N}\prod_{f\in\Fs}
\left(1+u^{m_s\deg(f)}\right).
\end{equation}
Thus $[u^{\kappa}]\Epoly(u,1)$ counts labeled codes of global code dimension
$\kappa$. This is a marginal distribution.
\end{corollary}

\begin{proof}
Each factor is either unselected, contributing its full global dimension weight,
or selected, contributing no code-dimension weight. The factor-selection bijection
and independence give the product.
\end{proof}

\begin{corollary}[Hull-dimension distribution]\label{cor:hull-dist}
The exact hull polynomial is
\begin{equation}\label{eq:hull-dist}
\mathcal H(z)=\Epoly(1,z)
 =\prod_{s=1}^{N}\prod_{O\in\Os}P_{a_O,w_O}(z).
\end{equation}
Consequently, $[z^{\eta}]\,\mathcal H(z)$ counts labeled codes with $k$-Galois
hull dimension $\eta$.
\end{corollary}

\begin{proof}
Set $u=1$ in Theorem~\ref{thm:central} and use
Propositions~\ref{prop:orbit-poly} and~\ref{prop:trace}.
\end{proof}

\begin{corollary}[LCD count]\label{cor:lcd}
All orbit weights are positive. Hence the hull dimension is zero exactly when
every orbit has $b_O=0$. A cyclic binary word has no $1$-to-$0$ transition exactly
when it is constant, so each orbit has two LCD selections. Therefore
\begin{equation}\label{eq:lcd}
\#\{C:\eta_k(C)=0\}=2^{\sum_s|\Os|}.
\end{equation}
This is a labeled LCD count.
\end{corollary}

\begin{proof}
The zero-boundary cyclic words are exactly the two constant words on each orbit.
Independence over the orbit partition gives the displayed product.
\end{proof}

For the moment statements, equip the finite labeled selection space
$\mathscr C$ with the uniform probability measure
\[
\mathbb P(C(J))=\frac{1}{\Epoly(1,1)}
 =2^{-\sum_s|\Fs|}
 \qquad(J\in\mathscr C).
\]
Thus every distinct labeled factor-selection code has equal probability.

\begin{corollary}[Mean hull dimension]\label{cor:mean}
Under this uniform measure, for an orbit of length $a\geq2$, the indicator
$X_i=\varepsilon_i(1-\varepsilon_{i+1})$ has expectation $1/4$, while a length-one
orbit contributes zero. Independence across orbits gives
\begin{equation}\label{eq:mean}
\mathbb E[\eta_k]
 =\sum_{s,O:\,a_O\geq2}\frac{a_Ow_O}{4}.
\end{equation}
Equivalently, the same value is obtained by differentiating the normalized hull
polynomial at $z=1$.
\end{corollary}

\begin{proof}
Sum the expectations of the directed boundary indicators over every orbit and
multiply by its weight. The length-one case has no directed boundary.
\end{proof}

\begin{corollary}[Variance of hull dimension]\label{cor:variance}
Under the same uniform measure, for $a\geq3$, neighboring directed-edge
indicators have covariance $-1/16$ and nonneighboring indicators have covariance
zero. Thus an orbit of length $a\geq3$ contributes $aw^2/16$. For $a=2$, the two
directed boundary indicators are mutually exclusive and the orbit contribution is
$w$ times a $\mathrm{Bernoulli}(1/2)$ variable, with variance $w^2/4$. A
length-one orbit contributes zero. Therefore
\begin{equation}\label{eq:var}
\operatorname{Var}(\eta_k)
 =\sum_{s,O:\,a_O=2}\frac{w_O^2}{4}
  +\sum_{s,O:\,a_O\geq3}\frac{a_Ow_O^2}{16}.
\end{equation}
The separate $a=2$ term is essential.
\end{corollary}

\begin{proof}
Under the uniform measure on $\mathscr C$ the indicators
$\varepsilon_{O,i}$ are independent $\mathrm{Bernoulli}(1/2)$ variables, because
$\mathscr C\cong\prod_{s,O}\{0,1\}^{a_O}$ (proof of
Theorem~\ref{thm:central}). Fix an orbit of length $a$ and weight $w$, and put
$X_i=\varepsilon_{O,i}(1-\varepsilon_{O,i+1})$ with indices modulo $a$. Its
contribution to $\eta_k$ is $w\sum_{i}X_i$.

\emph{Variances.} For an orbit of length at least two, $X_i$ is an indicator with
$\Pr(X_i=1)=1/4$, so $\operatorname{Var}X_i=\tfrac14-\tfrac1{16}=\tfrac3{16}$.

\emph{Case $a\ge3$.} Adjacent indicators satisfy
$X_iX_{i+1}=\varepsilon_{O,i}(1-\varepsilon_{O,i+1})\varepsilon_{O,i+1}(1-\varepsilon_{O,i+2})=0$,
hence $\operatorname{Cov}(X_i,X_{i+1})=0-\tfrac1{16}=-\tfrac1{16}$. If $X_i$ and
$X_j$ are not adjacent, i.e.\ $\{i,i+1\}\cap\{j,j+1\}=\emptyset$, they are
functions of disjoint sets of independent variables and hence independent. (For
$a=3$ every two distinct indicators are adjacent, so this case is empty.) Since
$a\ge3$, each $i$ has exactly two distinct neighbors $i\pm1$, so there are $2a$
ordered adjacent pairs. Therefore
\[
\operatorname{Var}\Bigl(\sum_i X_i\Bigr)
 = a\cdot\tfrac3{16}+2a\cdot\bigl(-\tfrac1{16}\bigr)
 = a\Bigl(\tfrac3{16}-\tfrac2{16}\Bigr)=\tfrac{a}{16},
\]
and the orbit contributes $a w^{2}/16$.

\emph{Case $a=2$.} Here $X_1=\varepsilon_{O,1}(1-\varepsilon_{O,2})$ and
$X_2=\varepsilon_{O,2}(1-\varepsilon_{O,1})$ are mutually exclusive, and
$X_1+X_2=1$ exactly when $\varepsilon_{O,1}\ne\varepsilon_{O,2}$. So $X_1+X_2$ is
$\mathrm{Bernoulli}(1/2)$ with variance $1/4$, and the weighted contribution is
$w^{2}/4$.

\emph{Case $a=1$.} $X_1=\varepsilon_{O,1}(1-\varepsilon_{O,1})=0$.

Distinct orbits involve disjoint independent variables, so their contributions are
independent and the variances add, giving \eqref{eq:var}.
\end{proof}

\begin{corollary}[Dimension-conditioned hull statistics]\label{cor:conditional}
For an integer $\kappa$ and $j\in\{0,1,2\}$ put
\[
\Phi_j(\kappa)=\sum_{\eta}\eta(\eta-1)\cdots(\eta-j+1)\,
 [u^{\kappa}z^{\eta}]\,\Epoly(u,z)
 =[u^{\kappa}]\,\frac{\partial^{j}\Epoly(u,z)}{\partial z^{j}}\Big|_{z=1},
\]
so that $\Phi_0(\kappa)=[u^{\kappa}]\,\Epoly(u,1)$ (empty product $=1$). Call
$\kappa$ attainable if $\Phi_0(\kappa)>0$. Under the uniform measure on labeled
selections of code dimension $\kappa$,
\begin{equation}\label{eq:conditional}
\Pr(\eta_k=\eta\mid\kappa(C)=\kappa)
 =\frac{[u^{\kappa}z^{\eta}]\,\Epoly(u,z)}{\Phi_0(\kappa)},\qquad
\mathbb E[\eta_k\mid\kappa(C)=\kappa]=\frac{\Phi_1(\kappa)}{\Phi_0(\kappa)}.
\end{equation}
These ratios are undefined when $\Phi_0(\kappa)=0$.
\end{corollary}

\begin{proof}
The denominator counts the $\Phi_0(\kappa)$ distinct labeled codes in the
conditioning event by Theorem~\ref{thm:central}. Its numerator counts those with
hull dimension $\eta$. Differentiating a monomial $z^{\eta}$ at $z=1$ weights it
by $\eta$, which gives $\Phi_1(\kappa)$ and proves the expectation formula.
\end{proof}

\begin{corollary}[Dimension-conditioned second moment and variance]\label{cor:conditional-variance}
Under the hypotheses and notation of Corollary~\ref{cor:conditional}, for
$\Phi_0(\kappa)>0$ the exact conditional factorial second moment and variance are
\begin{equation}\label{eq:conditional-variance}
\begin{gathered}
\mathbb E[\eta_k(\eta_k-1)\mid\kappa(C)=\kappa]=\frac{\Phi_2(\kappa)}{\Phi_0(\kappa)},
\\
\operatorname{Var}(\eta_k\mid\kappa(C)=\kappa)
 =\frac{\Phi_2(\kappa)+\Phi_1(\kappa)}{\Phi_0(\kappa)}
 -\Bigl(\frac{\Phi_1(\kappa)}{\Phi_0(\kappa)}\Bigr)^{2}.
\end{gathered}
\end{equation}
In particular these expressions are undefined if $\Phi_0(\kappa)=0$.
\end{corollary}

\begin{proof}
For each coefficient $c_{\kappa,\eta}=[u^{\kappa}z^{\eta}]\,\Epoly(u,z)$, first
and second $z$-differen\-tiation at $z=1$ give $\eta\,c_{\kappa,\eta}$ and
$\eta(\eta-1)\,c_{\kappa,\eta}$, respectively; summing over $\eta$ gives
$\Phi_1(\kappa)$ and $\Phi_2(\kappa)$. Divide by $\Phi_0(\kappa)$. Finally
$\eta^{2}=\eta(\eta-1)+\eta$ and the usual variance identity give
\eqref{eq:conditional-variance}. No independence of $\kappa$ and $\eta$ is
assumed.
\end{proof}

\subsection{Exact coefficient computation}\label{sec:algorithm}

\begin{proposition}[Orbit dynamic-programming cost]\label{prop:algorithm}
Assume the compatible orbit lengths and weights have already been computed. Put
$N_F=\sum_{s,O}a_O=\sum_s|\Fs|$ and $W=\sum_{s,O}a_Ow_O=n\sum_s m_s$. The full
bivariate coefficient polynomial \eqref{eq:joint} can be formed by tracking two
starting and current states around each orbit. A straightforward dense
implementation uses $O\bigl(N_F\,(W+1)(\lfloor W/2\rfloor+1)\bigr)$ integer
additions and $O\bigl((W+1)(\lfloor W/2\rfloor+1)\bigr)$ stored integer
coefficients. These are arithmetic-operation bounds, not bit-complexity bounds;
factorization and orbit construction are excluded.
\end{proposition}

\begin{proof}
\emph{Size of the array.} Code dimensions lie in $[0,W]$, and a nonconstant word
on an orbit of length $a$ has at most $\lfloor a/2\rfloor$ one-runs
(Proposition~\ref{prop:boundary-interpretation}), while a constant word
contributes zero to the hull dimension. Hence
$\eta\le\sum_{s,O}w_O\lfloor a_O/2\rfloor\le W/2$, and the coefficient array is
indexed by the rectangle $[0,W]\times[0,\lfloor W/2\rfloor]$. Intermediate open
walks may exceed this bound in the hull degree, but every later transition adds
nonnegative exponents, so a term leaving the rectangle cannot return to it; the
same rectangle therefore suffices at every step.

\emph{Recursion.} Keep a running global array $G$, initially the constant
polynomial $1$. For an orbit of length $a$, each starting state $s_0\in\{0,1\}$
initializes the current-state arrays with $G$ in state $s_0$ and $0$ in the other
state; a transition $(r,t)$ shifts an entry by $(w(1-t),\,wr(1-t))$ and adds it to
state $t$; after $a$ transitions the array in state $s_0$ collects the closed
walks. Summing the two starting states gives $G\cdot\operatorname{tr}(T_w^{a})$ by
Proposition~\ref{prop:trace}, and no separate convolution per orbit is needed.

\emph{Cost.} Each step performs four transitions for each of the two starting
states, each costing one pass over the rectangle, and there are $N_F$ steps;
a constant number of arrays gives the space bound. Coefficients can have
$O(N_F)$ bits, so integer additions are not unit-cost in a bit-complexity model.
\end{proof}

The coefficient calculation uses arbitrary-precision integers, and the
conditional moments are reduced rationals, so neither the integer coefficients
nor the reported fractional moments require floating-point rounding. Naive
enumeration visits all $2^{N_F}$ labeled selections, which the orbit recurrence
avoids; \S\ref{sec:examples} carries this out for a family with weight sum $W=31$.

\section{Illustrations and computational validation}\label{sec:examples}

We first apply Theorem~\ref{thm:central} to four families of increasing size,
including one whose orbit data already exceeds exhaustive enumeration, and then
report the computational validation.

\subsection{Worked examples}\label{sec:worked}

\begin{example}[Four \texorpdfstring{$\F_4$}{F4} components]\label{ex:pilot}
Consider the case $q=4$, $n=5$, $\lambda=1$, $k=1$ with four labeled $\F_4$
components, giving $8^4=4096$ labeled selections. Writing $\F_4=\F_2(a)$, the
factorization is
\[
x^5-1=(x+1)(x^2+ax+1)(x^2+(a+1)x+1),
\]
and the reciprocal action fixes the linear factor and interchanges the two
quadratics. Each component hence has orbits of lengths $1$ and $2$ with weights
$w=1$ and $w=2$. The one-component transfer-matrix traces of
Proposition~\ref{prop:trace} give the joint enumerator
\[
(u+1)(1+u^4+2u^2z^2)=1+u+u^4+u^5+2u^2z^2+2u^3z^2,
\]
which is the component joint histogram
$\{(0,0):1,(1,0):1,(2,2):2,(3,2):2,(4,0):1,(5,0):1\}$. The hull polynomial of one
component is $2(2+2z^2)=4+4z^2$, so the four-component ring hull polynomial is
\[
(4+4z^2)^4=256+1024z^2+1536z^4+1024z^6+256z^8,
\]
which is the ring hull histogram $\{0:256,2:1024,4:1536,6:1024,8:256\}$; the
ring joint histogram has $65$ terms. In particular, each component contributes
$4$ LCD selections and the ring contributes $256$, as given by
Corollary~\ref{cor:lcd}.
\end{example}

\begin{example}[A long orbit: \texorpdfstring{$\F_8$}{F8}, length \texorpdfstring{$7$}{7}]\label{ex:long-orbit}
Consider $q=8$, $n=7$, $\lambda=1$, $k=1$, with $2^7=128$ labeled selections by
Corollary~\ref{cor:total}. Here $x^7-1$ splits into seven linear factors over
$\F_8$, and the reciprocal action has orbit lengths $[1,6]$, all of weight $w=1$.
By \eqref{eq:orbit-poly}, the hull polynomial \eqref{eq:hull-dist} is therefore
\[
2P_{6,1}(z)=2(2+30z+30z^2+2z^3)=4+60z+60z^2+4z^3,
\]
which is the hull histogram $\{0:4,1:60,2:60,3:4\}$. The joint $(\kappa,\eta)$
multiplicities, in agreement with the transfer-matrix product of
Proposition~\ref{prop:trace}, are listed in Table~\ref{tab:F8}. The
$4$ LCD selections agree with Corollary~\ref{cor:lcd}. Among the
$\Phi_0(3)=35$ three-dimensional codes the table gives $12$, $21$ and $2$ with
hull dimensions $1$, $2$ and $3$, so the conditional mean at $\kappa=3$ is
$\Phi_1(3)/\Phi_0(3)=60/35=12/7$ by Corollary~\ref{cor:conditional}, and with
$\Phi_2(3)=54$, Corollary~\ref{cor:conditional-variance} gives the conditional
variance $78/245$.
\end{example}

\begin{table}[ht]
\centering
\caption{Joint multiplicities $[u^{\kappa}z^{\eta}]\,\Epoly(u,z)$ for $q=8$,
$n=7$, $\lambda=1$, $k=1$; a dash denotes a zero coefficient. Row $\kappa$ sums
to $\binom{7}{\kappa}$ and the table sums to $2^7=128$.}
\label{tab:F8}
\small
\begin{tabular}{@{}lrrrr@{}}
\toprule
$\kappa\backslash\eta$ & $0$ & $1$ & $2$ & $3$ \\
\midrule
$0$ & $1$ & -- & -- & -- \\
$1$ & $1$ & $6$ & -- & -- \\
$2$ & -- & $12$ & $9$ & -- \\
$3$ & -- & $12$ & $21$ & $2$ \\
$4$ & -- & $12$ & $21$ & $2$ \\
$5$ & -- & $12$ & $9$ & -- \\
$6$ & $1$ & $6$ & -- & -- \\
$7$ & $1$ & -- & -- & -- \\
\bottomrule
\end{tabular}
\end{table}

\begin{example}[A mixed algebra with extension weights]\label{ex:mixed}
Let $\A=\F_4\times\F_{16}$, $n=5$, $k=1$ and $\lambda=(\omega,1)$, where $\omega$
generates $\F_4^*$; then $256$ labeled selections are counted. The factors in
the $\F_4$ component give the orbits $(1,1)$ and $(2,2)$, while the five linear
factors over $\F_{16}$ give orbits $(1,2)$ and $(4,2)$: the weight of a factor in
the second component is $m_sd_O=2$, since $[\F_{16}:\F_4]=2$. The joint histogram
has $40$ terms and the hull polynomial is
\[
16+112z^2+112z^4+16z^6,
\]
of mean $3$ and variance $2$; the minimum $\eta_k=0$ is attained by $16=2^4$
codes, in agreement with Corollary~\ref{cor:lcd} since there are four orbits. The
example shows why the extension weights are essential: without the factor
$m_s=2$ the hull dimensions of the second component would be halved.
\end{example}

\begin{example}[A family beyond enumeration: \texorpdfstring{$\F_{32}$}{F32}, length \texorpdfstring{$31$}{31}]\label{ex:F32}
Let $q=32$, $n=31$, $\lambda=1$, $k=1$, so that $N_F=31$ and $W=31$. Naive
enumeration would visit all $2^{31}$ labeled selections. Write the roots of
$x^{31}-1$ as $\alpha^j$ with $\alpha$ primitive in $\F_{32}$ and
$j\in\mathbb Z/31\mathbb Z$; by Lemma~\ref{lem:root-action} the reciprocal acts
by $j\mapsto-16j\pmod{31}$, which has one fixed point and three cycles of length
$10$. Hence
\[
\Epoly(u,z)=(1+u)\bigl(\operatorname{tr}(T_1(u,z)^{10})\bigr)^3 .
\]
The exact expansion of \eqref{eq:joint} has $248$ nonzero coefficients, which sum
to $2^{31}$ and are computed by the recursion of \S\ref{sec:algorithm} without
enumerating selections; the whole distribution is displayed in
Fig.~\ref{fig:joint}. Its support is the diamond
$0\le\eta\le\min(\kappa,31-\kappa)$ predicted by nondegeneracy, and it is
symmetric under $\kappa\mapsto31-\kappa$. The enumerator gives $16=2^4$ LCD codes,
while Corollaries~\ref{cor:mean} and~\ref{cor:variance} give the exact mean
$15/2$ and variance $15/8$ of the hull dimension. The conditional mean
$\mathbb E[\eta_k\mid\kappa]$ rises to about $7.74$ at $\kappa=15$ and decreases
symmetrically; for instance $\mathbb E[\eta_k\mid\kappa=1]=30/31$ and
$\mathbb E[\eta_k\mid\kappa=15]=240/31$.
\end{example}

\begin{figure}[htbp]
\centering
\begin{tikzpicture}[x=1mm,y=1mm]
  \fill[black!18] (37.5,37.5) rectangle +(2.5,2.5);
  \fill[black!18] (40.0,37.5) rectangle +(2.5,2.5);
  \fill[black!37] (35.0,35.0) rectangle +(2.5,2.5);
  \fill[black!42] (37.5,35.0) rectangle +(2.5,2.5);
  \fill[black!42] (40.0,35.0) rectangle +(2.5,2.5);
  \fill[black!37] (42.5,35.0) rectangle +(2.5,2.5);
  \fill[black!50] (32.5,32.5) rectangle +(2.5,2.5);
  \fill[black!57] (35.0,32.5) rectangle +(2.5,2.5);
  \fill[black!60] (37.5,32.5) rectangle +(2.5,2.5);
  \fill[black!60] (40.0,32.5) rectangle +(2.5,2.5);
  \fill[black!57] (42.5,32.5) rectangle +(2.5,2.5);
  \fill[black!50] (45.0,32.5) rectangle +(2.5,2.5);
  \fill[black!60] (30.0,30.0) rectangle +(2.5,2.5);
  \fill[black!68] (32.5,30.0) rectangle +(2.5,2.5);
  \fill[black!73] (35.0,30.0) rectangle +(2.5,2.5);
  \fill[black!75] (37.5,30.0) rectangle +(2.5,2.5);
  \fill[black!75] (40.0,30.0) rectangle +(2.5,2.5);
  \fill[black!73] (42.5,30.0) rectangle +(2.5,2.5);
  \fill[black!68] (45.0,30.0) rectangle +(2.5,2.5);
  \fill[black!60] (47.5,30.0) rectangle +(2.5,2.5);
  \fill[black!67] (27.5,27.5) rectangle +(2.5,2.5);
  \fill[black!76] (30.0,27.5) rectangle +(2.5,2.5);
  \fill[black!81] (32.5,27.5) rectangle +(2.5,2.5);
  \fill[black!84] (35.0,27.5) rectangle +(2.5,2.5);
  \fill[black!86] (37.5,27.5) rectangle +(2.5,2.5);
  \fill[black!86] (40.0,27.5) rectangle +(2.5,2.5);
  \fill[black!84] (42.5,27.5) rectangle +(2.5,2.5);
  \fill[black!81] (45.0,27.5) rectangle +(2.5,2.5);
  \fill[black!76] (47.5,27.5) rectangle +(2.5,2.5);
  \fill[black!67] (50.0,27.5) rectangle +(2.5,2.5);
  \fill[black!71] (25.0,25.0) rectangle +(2.5,2.5);
  \fill[black!80] (27.5,25.0) rectangle +(2.5,2.5);
  \fill[black!86] (30.0,25.0) rectangle +(2.5,2.5);
  \fill[black!90] (32.5,25.0) rectangle +(2.5,2.5);
  \fill[black!92] (35.0,25.0) rectangle +(2.5,2.5);
  \fill[black!93] (37.5,25.0) rectangle +(2.5,2.5);
  \fill[black!93] (40.0,25.0) rectangle +(2.5,2.5);
  \fill[black!92] (42.5,25.0) rectangle +(2.5,2.5);
  \fill[black!90] (45.0,25.0) rectangle +(2.5,2.5);
  \fill[black!86] (47.5,25.0) rectangle +(2.5,2.5);
  \fill[black!80] (50.0,25.0) rectangle +(2.5,2.5);
  \fill[black!71] (52.5,25.0) rectangle +(2.5,2.5);
  \fill[black!73] (22.5,22.5) rectangle +(2.5,2.5);
  \fill[black!83] (25.0,22.5) rectangle +(2.5,2.5);
  \fill[black!89] (27.5,22.5) rectangle +(2.5,2.5);
  \fill[black!93] (30.0,22.5) rectangle +(2.5,2.5);
  \fill[black!95] (32.5,22.5) rectangle +(2.5,2.5);
  \fill[black!97] (35.0,22.5) rectangle +(2.5,2.5);
  \fill[black!98] (37.5,22.5) rectangle +(2.5,2.5);
  \fill[black!98] (40.0,22.5) rectangle +(2.5,2.5);
  \fill[black!97] (42.5,22.5) rectangle +(2.5,2.5);
  \fill[black!95] (45.0,22.5) rectangle +(2.5,2.5);
  \fill[black!93] (47.5,22.5) rectangle +(2.5,2.5);
  \fill[black!89] (50.0,22.5) rectangle +(2.5,2.5);
  \fill[black!83] (52.5,22.5) rectangle +(2.5,2.5);
  \fill[black!73] (55.0,22.5) rectangle +(2.5,2.5);
  \fill[black!73] (20.0,20.0) rectangle +(2.5,2.5);
  \fill[black!83] (22.5,20.0) rectangle +(2.5,2.5);
  \fill[black!89] (25.0,20.0) rectangle +(2.5,2.5);
  \fill[black!93] (27.5,20.0) rectangle +(2.5,2.5);
  \fill[black!96] (30.0,20.0) rectangle +(2.5,2.5);
  \fill[black!98] (32.5,20.0) rectangle +(2.5,2.5);
  \fill[black!99] (35.0,20.0) rectangle +(2.5,2.5);
  \fill[black!100] (37.5,20.0) rectangle +(2.5,2.5);
  \fill[black!100] (40.0,20.0) rectangle +(2.5,2.5);
  \fill[black!99] (42.5,20.0) rectangle +(2.5,2.5);
  \fill[black!98] (45.0,20.0) rectangle +(2.5,2.5);
  \fill[black!96] (47.5,20.0) rectangle +(2.5,2.5);
  \fill[black!93] (50.0,20.0) rectangle +(2.5,2.5);
  \fill[black!89] (52.5,20.0) rectangle +(2.5,2.5);
  \fill[black!83] (55.0,20.0) rectangle +(2.5,2.5);
  \fill[black!73] (57.5,20.0) rectangle +(2.5,2.5);
  \fill[black!72] (17.5,17.5) rectangle +(2.5,2.5);
  \fill[black!81] (20.0,17.5) rectangle +(2.5,2.5);
  \fill[black!87] (22.5,17.5) rectangle +(2.5,2.5);
  \fill[black!91] (25.0,17.5) rectangle +(2.5,2.5);
  \fill[black!94] (27.5,17.5) rectangle +(2.5,2.5);
  \fill[black!96] (30.0,17.5) rectangle +(2.5,2.5);
  \fill[black!98] (32.5,17.5) rectangle +(2.5,2.5);
  \fill[black!99] (35.0,17.5) rectangle +(2.5,2.5);
  \fill[black!99] (37.5,17.5) rectangle +(2.5,2.5);
  \fill[black!99] (40.0,17.5) rectangle +(2.5,2.5);
  \fill[black!99] (42.5,17.5) rectangle +(2.5,2.5);
  \fill[black!98] (45.0,17.5) rectangle +(2.5,2.5);
  \fill[black!96] (47.5,17.5) rectangle +(2.5,2.5);
  \fill[black!94] (50.0,17.5) rectangle +(2.5,2.5);
  \fill[black!91] (52.5,17.5) rectangle +(2.5,2.5);
  \fill[black!87] (55.0,17.5) rectangle +(2.5,2.5);
  \fill[black!81] (57.5,17.5) rectangle +(2.5,2.5);
  \fill[black!72] (60.0,17.5) rectangle +(2.5,2.5);
  \fill[black!69] (15.0,15.0) rectangle +(2.5,2.5);
  \fill[black!77] (17.5,15.0) rectangle +(2.5,2.5);
  \fill[black!83] (20.0,15.0) rectangle +(2.5,2.5);
  \fill[black!87] (22.5,15.0) rectangle +(2.5,2.5);
  \fill[black!90] (25.0,15.0) rectangle +(2.5,2.5);
  \fill[black!92] (27.5,15.0) rectangle +(2.5,2.5);
  \fill[black!94] (30.0,15.0) rectangle +(2.5,2.5);
  \fill[black!95] (32.5,15.0) rectangle +(2.5,2.5);
  \fill[black!96] (35.0,15.0) rectangle +(2.5,2.5);
  \fill[black!96] (37.5,15.0) rectangle +(2.5,2.5);
  \fill[black!96] (40.0,15.0) rectangle +(2.5,2.5);
  \fill[black!96] (42.5,15.0) rectangle +(2.5,2.5);
  \fill[black!95] (45.0,15.0) rectangle +(2.5,2.5);
  \fill[black!94] (47.5,15.0) rectangle +(2.5,2.5);
  \fill[black!92] (50.0,15.0) rectangle +(2.5,2.5);
  \fill[black!90] (52.5,15.0) rectangle +(2.5,2.5);
  \fill[black!87] (55.0,15.0) rectangle +(2.5,2.5);
  \fill[black!83] (57.5,15.0) rectangle +(2.5,2.5);
  \fill[black!77] (60.0,15.0) rectangle +(2.5,2.5);
  \fill[black!69] (62.5,15.0) rectangle +(2.5,2.5);
  \fill[black!64] (12.5,12.5) rectangle +(2.5,2.5);
  \fill[black!72] (15.0,12.5) rectangle +(2.5,2.5);
  \fill[black!77] (17.5,12.5) rectangle +(2.5,2.5);
  \fill[black!81] (20.0,12.5) rectangle +(2.5,2.5);
  \fill[black!84] (22.5,12.5) rectangle +(2.5,2.5);
  \fill[black!86] (25.0,12.5) rectangle +(2.5,2.5);
  \fill[black!87] (27.5,12.5) rectangle +(2.5,2.5);
  \fill[black!88] (30.0,12.5) rectangle +(2.5,2.5);
  \fill[black!89] (32.5,12.5) rectangle +(2.5,2.5);
  \fill[black!90] (35.0,12.5) rectangle +(2.5,2.5);
  \fill[black!90] (37.5,12.5) rectangle +(2.5,2.5);
  \fill[black!90] (40.0,12.5) rectangle +(2.5,2.5);
  \fill[black!90] (42.5,12.5) rectangle +(2.5,2.5);
  \fill[black!89] (45.0,12.5) rectangle +(2.5,2.5);
  \fill[black!88] (47.5,12.5) rectangle +(2.5,2.5);
  \fill[black!87] (50.0,12.5) rectangle +(2.5,2.5);
  \fill[black!86] (52.5,12.5) rectangle +(2.5,2.5);
  \fill[black!84] (55.0,12.5) rectangle +(2.5,2.5);
  \fill[black!81] (57.5,12.5) rectangle +(2.5,2.5);
  \fill[black!77] (60.0,12.5) rectangle +(2.5,2.5);
  \fill[black!72] (62.5,12.5) rectangle +(2.5,2.5);
  \fill[black!64] (65.0,12.5) rectangle +(2.5,2.5);
  \fill[black!57] (10.0,10.0) rectangle +(2.5,2.5);
  \fill[black!65] (12.5,10.0) rectangle +(2.5,2.5);
  \fill[black!69] (15.0,10.0) rectangle +(2.5,2.5);
  \fill[black!72] (17.5,10.0) rectangle +(2.5,2.5);
  \fill[black!75] (20.0,10.0) rectangle +(2.5,2.5);
  \fill[black!77] (22.5,10.0) rectangle +(2.5,2.5);
  \fill[black!78] (25.0,10.0) rectangle +(2.5,2.5);
  \fill[black!79] (27.5,10.0) rectangle +(2.5,2.5);
  \fill[black!80] (30.0,10.0) rectangle +(2.5,2.5);
  \fill[black!80] (32.5,10.0) rectangle +(2.5,2.5);
  \fill[black!81] (35.0,10.0) rectangle +(2.5,2.5);
  \fill[black!81] (37.5,10.0) rectangle +(2.5,2.5);
  \fill[black!81] (40.0,10.0) rectangle +(2.5,2.5);
  \fill[black!81] (42.5,10.0) rectangle +(2.5,2.5);
  \fill[black!80] (45.0,10.0) rectangle +(2.5,2.5);
  \fill[black!80] (47.5,10.0) rectangle +(2.5,2.5);
  \fill[black!79] (50.0,10.0) rectangle +(2.5,2.5);
  \fill[black!78] (52.5,10.0) rectangle +(2.5,2.5);
  \fill[black!77] (55.0,10.0) rectangle +(2.5,2.5);
  \fill[black!75] (57.5,10.0) rectangle +(2.5,2.5);
  \fill[black!72] (60.0,10.0) rectangle +(2.5,2.5);
  \fill[black!69] (62.5,10.0) rectangle +(2.5,2.5);
  \fill[black!65] (65.0,10.0) rectangle +(2.5,2.5);
  \fill[black!57] (67.5,10.0) rectangle +(2.5,2.5);
  \fill[black!49] (7.5,7.5) rectangle +(2.5,2.5);
  \fill[black!55] (10.0,7.5) rectangle +(2.5,2.5);
  \fill[black!59] (12.5,7.5) rectangle +(2.5,2.5);
  \fill[black!62] (15.0,7.5) rectangle +(2.5,2.5);
  \fill[black!63] (17.5,7.5) rectangle +(2.5,2.5);
  \fill[black!65] (20.0,7.5) rectangle +(2.5,2.5);
  \fill[black!66] (22.5,7.5) rectangle +(2.5,2.5);
  \fill[black!67] (25.0,7.5) rectangle +(2.5,2.5);
  \fill[black!67] (27.5,7.5) rectangle +(2.5,2.5);
  \fill[black!68] (30.0,7.5) rectangle +(2.5,2.5);
  \fill[black!68] (32.5,7.5) rectangle +(2.5,2.5);
  \fill[black!68] (35.0,7.5) rectangle +(2.5,2.5);
  \fill[black!68] (37.5,7.5) rectangle +(2.5,2.5);
  \fill[black!68] (40.0,7.5) rectangle +(2.5,2.5);
  \fill[black!68] (42.5,7.5) rectangle +(2.5,2.5);
  \fill[black!68] (45.0,7.5) rectangle +(2.5,2.5);
  \fill[black!68] (47.5,7.5) rectangle +(2.5,2.5);
  \fill[black!67] (50.0,7.5) rectangle +(2.5,2.5);
  \fill[black!67] (52.5,7.5) rectangle +(2.5,2.5);
  \fill[black!66] (55.0,7.5) rectangle +(2.5,2.5);
  \fill[black!65] (57.5,7.5) rectangle +(2.5,2.5);
  \fill[black!63] (60.0,7.5) rectangle +(2.5,2.5);
  \fill[black!62] (62.5,7.5) rectangle +(2.5,2.5);
  \fill[black!59] (65.0,7.5) rectangle +(2.5,2.5);
  \fill[black!55] (67.5,7.5) rectangle +(2.5,2.5);
  \fill[black!49] (70.0,7.5) rectangle +(2.5,2.5);
  \fill[black!38] (5.0,5.0) rectangle +(2.5,2.5);
  \fill[black!44] (7.5,5.0) rectangle +(2.5,2.5);
  \fill[black!46] (10.0,5.0) rectangle +(2.5,2.5);
  \fill[black!48] (12.5,5.0) rectangle +(2.5,2.5);
  \fill[black!49] (15.0,5.0) rectangle +(2.5,2.5);
  \fill[black!49] (17.5,5.0) rectangle +(2.5,2.5);
  \fill[black!50] (20.0,5.0) rectangle +(2.5,2.5);
  \fill[black!50] (22.5,5.0) rectangle +(2.5,2.5);
  \fill[black!51] (25.0,5.0) rectangle +(2.5,2.5);
  \fill[black!51] (27.5,5.0) rectangle +(2.5,2.5);
  \fill[black!51] (30.0,5.0) rectangle +(2.5,2.5);
  \fill[black!51] (32.5,5.0) rectangle +(2.5,2.5);
  \fill[black!51] (35.0,5.0) rectangle +(2.5,2.5);
  \fill[black!52] (37.5,5.0) rectangle +(2.5,2.5);
  \fill[black!52] (40.0,5.0) rectangle +(2.5,2.5);
  \fill[black!51] (42.5,5.0) rectangle +(2.5,2.5);
  \fill[black!51] (45.0,5.0) rectangle +(2.5,2.5);
  \fill[black!51] (47.5,5.0) rectangle +(2.5,2.5);
  \fill[black!51] (50.0,5.0) rectangle +(2.5,2.5);
  \fill[black!51] (52.5,5.0) rectangle +(2.5,2.5);
  \fill[black!50] (55.0,5.0) rectangle +(2.5,2.5);
  \fill[black!50] (57.5,5.0) rectangle +(2.5,2.5);
  \fill[black!49] (60.0,5.0) rectangle +(2.5,2.5);
  \fill[black!49] (62.5,5.0) rectangle +(2.5,2.5);
  \fill[black!48] (65.0,5.0) rectangle +(2.5,2.5);
  \fill[black!46] (67.5,5.0) rectangle +(2.5,2.5);
  \fill[black!44] (70.0,5.0) rectangle +(2.5,2.5);
  \fill[black!38] (72.5,5.0) rectangle +(2.5,2.5);
  \fill[black!25] (2.5,2.5) rectangle +(2.5,2.5);
  \fill[black!29] (5.0,2.5) rectangle +(2.5,2.5);
  \fill[black!29] (7.5,2.5) rectangle +(2.5,2.5);
  \fill[black!29] (10.0,2.5) rectangle +(2.5,2.5);
  \fill[black!29] (12.5,2.5) rectangle +(2.5,2.5);
  \fill[black!29] (15.0,2.5) rectangle +(2.5,2.5);
  \fill[black!29] (17.5,2.5) rectangle +(2.5,2.5);
  \fill[black!29] (20.0,2.5) rectangle +(2.5,2.5);
  \fill[black!29] (22.5,2.5) rectangle +(2.5,2.5);
  \fill[black!25] (25.0,2.5) rectangle +(2.5,2.5);
  \fill[black!29] (27.5,2.5) rectangle +(2.5,2.5);
  \fill[black!32] (30.0,2.5) rectangle +(2.5,2.5);
  \fill[black!32] (32.5,2.5) rectangle +(2.5,2.5);
  \fill[black!32] (35.0,2.5) rectangle +(2.5,2.5);
  \fill[black!32] (37.5,2.5) rectangle +(2.5,2.5);
  \fill[black!32] (40.0,2.5) rectangle +(2.5,2.5);
  \fill[black!32] (42.5,2.5) rectangle +(2.5,2.5);
  \fill[black!32] (45.0,2.5) rectangle +(2.5,2.5);
  \fill[black!32] (47.5,2.5) rectangle +(2.5,2.5);
  \fill[black!29] (50.0,2.5) rectangle +(2.5,2.5);
  \fill[black!25] (52.5,2.5) rectangle +(2.5,2.5);
  \fill[black!29] (55.0,2.5) rectangle +(2.5,2.5);
  \fill[black!29] (57.5,2.5) rectangle +(2.5,2.5);
  \fill[black!29] (60.0,2.5) rectangle +(2.5,2.5);
  \fill[black!29] (62.5,2.5) rectangle +(2.5,2.5);
  \fill[black!29] (65.0,2.5) rectangle +(2.5,2.5);
  \fill[black!29] (67.5,2.5) rectangle +(2.5,2.5);
  \fill[black!29] (70.0,2.5) rectangle +(2.5,2.5);
  \fill[black!29] (72.5,2.5) rectangle +(2.5,2.5);
  \fill[black!25] (75.0,2.5) rectangle +(2.5,2.5);
  \fill[black!8] (0.0,0.0) rectangle +(2.5,2.5);
  \fill[black!8] (2.5,0.0) rectangle +(2.5,2.5);
  \fill[black!14] (25.0,0.0) rectangle +(2.5,2.5);
  \fill[black!14] (27.5,0.0) rectangle +(2.5,2.5);
  \fill[black!14] (50.0,0.0) rectangle +(2.5,2.5);
  \fill[black!14] (52.5,0.0) rectangle +(2.5,2.5);
  \fill[black!8] (75.0,0.0) rectangle +(2.5,2.5);
  \fill[black!8] (77.5,0.0) rectangle +(2.5,2.5);
  \draw[black!30,line width=0.15pt] (0,0) grid[step=2.5] (80.0,40.0);
  \draw[black,line width=0.4pt] (0,0) rectangle (80.0,40.0);
  \draw[black,line width=0.4pt] (1.25,0) -- +(0,-1.2mm);
  \node[font=\scriptsize,anchor=north] at (1.25,-1.8) {0};
  \draw[black,line width=0.4pt] (13.75,0) -- +(0,-1.2mm);
  \node[font=\scriptsize,anchor=north] at (13.75,-1.8) {5};
  \draw[black,line width=0.4pt] (26.25,0) -- +(0,-1.2mm);
  \node[font=\scriptsize,anchor=north] at (26.25,-1.8) {10};
  \draw[black,line width=0.4pt] (38.75,0) -- +(0,-1.2mm);
  \node[font=\scriptsize,anchor=north] at (38.75,-1.8) {15};
  \draw[black,line width=0.4pt] (51.25,0) -- +(0,-1.2mm);
  \node[font=\scriptsize,anchor=north] at (51.25,-1.8) {20};
  \draw[black,line width=0.4pt] (63.75,0) -- +(0,-1.2mm);
  \node[font=\scriptsize,anchor=north] at (63.75,-1.8) {25};
  \draw[black,line width=0.4pt] (76.25,0) -- +(0,-1.2mm);
  \node[font=\scriptsize,anchor=north] at (76.25,-1.8) {30};
  \draw[black,line width=0.4pt] (0,1.25) -- +(-1.2mm,0);
  \node[font=\scriptsize,anchor=east] at (-1.8,1.25) {0};
  \draw[black,line width=0.4pt] (0,13.75) -- +(-1.2mm,0);
  \node[font=\scriptsize,anchor=east] at (-1.8,13.75) {5};
  \draw[black,line width=0.4pt] (0,26.25) -- +(-1.2mm,0);
  \node[font=\scriptsize,anchor=east] at (-1.8,26.25) {10};
  \draw[black,line width=0.4pt] (0,38.75) -- +(-1.2mm,0);
  \node[font=\scriptsize,anchor=east] at (-1.8,38.75) {15};
  \draw[black,line width=0.6pt] (1.25,1.25) -- (3.75,3.67) -- (6.25,5.93) -- (8.75,8.02) -- (11.25,9.96) -- (13.75,11.73) -- (16.25,13.35) -- (18.75,14.80) -- (21.25,16.09) -- (23.75,17.22) -- (26.25,18.19) -- (28.75,18.99) -- (31.25,19.64) -- (33.75,20.12) -- (36.25,20.44) -- (38.75,20.60) -- (41.25,20.60) -- (43.75,20.44) -- (46.25,20.12) -- (48.75,19.64) -- (51.25,18.99) -- (53.75,18.19) -- (56.25,17.22) -- (58.75,16.09) -- (61.25,14.80) -- (63.75,13.35) -- (66.25,11.73) -- (68.75,9.96) -- (71.25,8.02) -- (73.75,5.93) -- (76.25,3.67) -- (78.75,1.25);
  \fill[white] (1.25,1.25) circle (0.5mm);
  \draw[black,line width=0.3pt] (1.25,1.25) circle (0.5mm);
  \fill[white] (3.75,3.67) circle (0.5mm);
  \draw[black,line width=0.3pt] (3.75,3.67) circle (0.5mm);
  \fill[white] (6.25,5.93) circle (0.5mm);
  \draw[black,line width=0.3pt] (6.25,5.93) circle (0.5mm);
  \fill[white] (8.75,8.02) circle (0.5mm);
  \draw[black,line width=0.3pt] (8.75,8.02) circle (0.5mm);
  \fill[white] (11.25,9.96) circle (0.5mm);
  \draw[black,line width=0.3pt] (11.25,9.96) circle (0.5mm);
  \fill[white] (13.75,11.73) circle (0.5mm);
  \draw[black,line width=0.3pt] (13.75,11.73) circle (0.5mm);
  \fill[white] (16.25,13.35) circle (0.5mm);
  \draw[black,line width=0.3pt] (16.25,13.35) circle (0.5mm);
  \fill[white] (18.75,14.80) circle (0.5mm);
  \draw[black,line width=0.3pt] (18.75,14.80) circle (0.5mm);
  \fill[white] (21.25,16.09) circle (0.5mm);
  \draw[black,line width=0.3pt] (21.25,16.09) circle (0.5mm);
  \fill[white] (23.75,17.22) circle (0.5mm);
  \draw[black,line width=0.3pt] (23.75,17.22) circle (0.5mm);
  \fill[white] (26.25,18.19) circle (0.5mm);
  \draw[black,line width=0.3pt] (26.25,18.19) circle (0.5mm);
  \fill[white] (28.75,18.99) circle (0.5mm);
  \draw[black,line width=0.3pt] (28.75,18.99) circle (0.5mm);
  \fill[white] (31.25,19.64) circle (0.5mm);
  \draw[black,line width=0.3pt] (31.25,19.64) circle (0.5mm);
  \fill[white] (33.75,20.12) circle (0.5mm);
  \draw[black,line width=0.3pt] (33.75,20.12) circle (0.5mm);
  \fill[white] (36.25,20.44) circle (0.5mm);
  \draw[black,line width=0.3pt] (36.25,20.44) circle (0.5mm);
  \fill[white] (38.75,20.60) circle (0.5mm);
  \draw[black,line width=0.3pt] (38.75,20.60) circle (0.5mm);
  \fill[white] (41.25,20.60) circle (0.5mm);
  \draw[black,line width=0.3pt] (41.25,20.60) circle (0.5mm);
  \fill[white] (43.75,20.44) circle (0.5mm);
  \draw[black,line width=0.3pt] (43.75,20.44) circle (0.5mm);
  \fill[white] (46.25,20.12) circle (0.5mm);
  \draw[black,line width=0.3pt] (46.25,20.12) circle (0.5mm);
  \fill[white] (48.75,19.64) circle (0.5mm);
  \draw[black,line width=0.3pt] (48.75,19.64) circle (0.5mm);
  \fill[white] (51.25,18.99) circle (0.5mm);
  \draw[black,line width=0.3pt] (51.25,18.99) circle (0.5mm);
  \fill[white] (53.75,18.19) circle (0.5mm);
  \draw[black,line width=0.3pt] (53.75,18.19) circle (0.5mm);
  \fill[white] (56.25,17.22) circle (0.5mm);
  \draw[black,line width=0.3pt] (56.25,17.22) circle (0.5mm);
  \fill[white] (58.75,16.09) circle (0.5mm);
  \draw[black,line width=0.3pt] (58.75,16.09) circle (0.5mm);
  \fill[white] (61.25,14.80) circle (0.5mm);
  \draw[black,line width=0.3pt] (61.25,14.80) circle (0.5mm);
  \fill[white] (63.75,13.35) circle (0.5mm);
  \draw[black,line width=0.3pt] (63.75,13.35) circle (0.5mm);
  \fill[white] (66.25,11.73) circle (0.5mm);
  \draw[black,line width=0.3pt] (66.25,11.73) circle (0.5mm);
  \fill[white] (68.75,9.96) circle (0.5mm);
  \draw[black,line width=0.3pt] (68.75,9.96) circle (0.5mm);
  \fill[white] (71.25,8.02) circle (0.5mm);
  \draw[black,line width=0.3pt] (71.25,8.02) circle (0.5mm);
  \fill[white] (73.75,5.93) circle (0.5mm);
  \draw[black,line width=0.3pt] (73.75,5.93) circle (0.5mm);
  \fill[white] (76.25,3.67) circle (0.5mm);
  \draw[black,line width=0.3pt] (76.25,3.67) circle (0.5mm);
  \fill[white] (78.75,1.25) circle (0.5mm);
  \draw[black,line width=0.3pt] (78.75,1.25) circle (0.5mm);
  \node[font=\small,anchor=north] at (40.0,-4.6) {$\kappa$};
  \node[font=\small,anchor=south] at (-7.5,20.0) {$\eta_k$};
  \fill[black!8] (0.0,-13.0) rectangle +(3mm,3mm);
  \fill[black!18] (3.0,-13.0) rectangle +(3mm,3mm);
  \fill[black!28] (6.0,-13.0) rectangle +(3mm,3mm);
  \fill[black!39] (9.0,-13.0) rectangle +(3mm,3mm);
  \fill[black!49] (12.0,-13.0) rectangle +(3mm,3mm);
  \fill[black!59] (15.0,-13.0) rectangle +(3mm,3mm);
  \fill[black!69] (18.0,-13.0) rectangle +(3mm,3mm);
  \fill[black!80] (21.0,-13.0) rectangle +(3mm,3mm);
  \fill[black!90] (24.0,-13.0) rectangle +(3mm,3mm);
  \fill[black!100] (27.0,-13.0) rectangle +(3mm,3mm);
  \draw[black,line width=0.4pt] (0,-13.0) rectangle (30.0,-10.0);
  \node[font=\scriptsize,anchor=west,align=left,text width=34mm] at (32.5,-11.5) {multiplicity, $\log$ scale};
  \node[font=\scriptsize,anchor=west,align=left,text width=40mm] at (32.5,-15.5) {$1$ (light) to $8.55\times10^{7}$ (dark)};
  \fill[white] (1.5,-18.0) circle (0.5mm);
  \draw[black,line width=0.3pt] (1.5,-18.0) circle (0.5mm);
  \node[font=\scriptsize,anchor=west] at (4.0,-18.0) {exact conditional mean $\mathbb E[\eta_k\mid\kappa]$};
\end{tikzpicture}
\caption{Exact joint multiplicities $[u^{\kappa}z^{\eta}]\,\Epoly(u,z)$ for
$\F_{32}$, $n=31$, $\lambda=1$, $k=1$, shaded on a logarithmic scale. The $248$
nonzero coefficients, obtained from the orbit data of Example~\ref{ex:F32} by the
recurrence of \S\ref{sec:algorithm}, account exactly for all $2^{31}$ labeled
factor selections. Nonzero coefficients fill the diamond
$0\le\eta\le\min(\kappa,31-\kappa)$ and are symmetric under
$\kappa\mapsto31-\kappa$; the white circles mark the exact conditional mean
$\mathbb E[\eta_k\mid\kappa]$.}
\label{fig:joint}
\end{figure}

\subsection{Computational validation}\label{sec:verification}

Each joint enumerator is evaluated twice: by the trace recursion of
Proposition~\ref{prop:algorithm} and through the closed-form coefficients of
Proposition~\ref{prop:closed} followed by convolution over the orbits. The two
evaluations agree coefficient by coefficient for every family of
Table~\ref{tab:validation}. All coefficient arithmetic is integer arithmetic on
the orbit data, and the conditional moments of Corollaries~\ref{cor:conditional}
and~\ref{cor:conditional-variance} are reduced rationals, so no floating-point
rounding enters the reported values.

Independently, hull dimensions of individual codes are computed from the
definitions rather than from the orbit formula. For a selection $J$ a generator
matrix $B$ of $C(J)$ is formed in coefficient coordinates, the $k$-Galois dual is
obtained as $\rho_{s,k}(\operatorname{null}B)$ over the component field, and the
hull dimension from
\[
\dim\Hull_k(C)=\dim C+\dim D-\dim(C+D).
\]
All ranks and null spaces are computed over the component field. In each of the
small families of Table~\ref{tab:validation} this reproduces every hull dimension
predicted by the orbit formula, and the joint histograms agree as well. The
families in characteristics $3$ and $5$ are included because signs and
normalizations in \eqref{eq:reciprocal} are invisible when $-1=1$, so they
exercise the formulas in a different arithmetic setting. For $\F_{32}$ and
$n=31$ individual codes agree with the orbit formula, while the enumerator itself
counts all $2^{31}$ selections exactly through the orbit recurrence.

For the mixed algebra $\A=\F_4[X]/\langle X(X^2+X+\omega)\rangle$ with
$\lambda=\omega(1+X+X^2)$ the same count was obtained directly in
polynomial-quotient coordinates. In the basis $1,X,X^2$ each $\A$-valued
orthogonality condition becomes three equations over $\F_4$. There the Frobenius
map is ring squaring, not coefficientwise squaring with $X$ fixed; since squaring
is bijective on $\F_4$, the nullity of the resulting linear system is the hull
dimension. The result agrees with the componentwise enumeration of
Example~\ref{ex:mixed}.

\begin{table}[ht]
\centering
\caption{Representative families and enumerative results. Dimensions of the
$\F_{16}$ component are taken over $\F_4$, so its orbit weights contain the
extension degree $[\F_{16}:\F_4]=2$; here $\omega$ generates $\F_4^*$ and
$i\in\F_9$ satisfies $i^2=-1$. The last column lists the independent check
described in \S\ref{sec:verification}, and ``codes'' is the number of labeled
factor selections.}
\label{tab:validation}
\footnotesize
\renewcommand{\arraystretch}{1.25}
\setlength{\tabcolsep}{3.5pt}
\begin{tabularx}{\textwidth}{@{}L{25mm}L{19mm}rL{28mm}>{\raggedright\arraybackslash}X@{}}
\toprule
Family and parameters & Reciprocal orbits $(a_O,w_O)$ & Codes & Principal exact result & Independent check \\
\midrule
$\F_4^{\times4}$\newline $n=5$, $\lambda=1$, $k=1$ & $(1,1)$, $(2,2)$ per component & $4096$ & $65$ joint terms; $256$ LCD & all selections \\
$\F_8$\newline $n=7$, $\lambda=1$, $k=1$ & $(1,1)$, $(6,1)$ & $128$ & $16$ joint terms (Table~\ref{tab:F8}); $\mathbb E[\eta_k\mid\kappa=3]=12/7$ & all codes \\
$\F_{16}\!/\!\F_4$\newline $n=5$, $\lambda=1$, $k=1$ & $(1,2)$, $(4,2)$ & $32$ & mean $2$, variance $1$ & all codes \\
$\F_4$\newline $n=5$, $\lambda=\omega$, $k=1$ & $(1,1)$, $(2,2)$ & $8$ & $4$ LCD & all codes \\
$\F_{4}\!\times\!\F_{16}$\newline $n=5$, $\lambda=(\omega,1)$, $k=1$ & $(1,1),(2,2)$, $(1,2),(4,2)$ & $256$ & $40$ joint terms; mean $3$, variance $2$ & all codes; quotient coordinates \\
$\F_3$\newline $n=4$, $\lambda=1,-1$, $k=0$ & --- & $8,4$ & distributions from the orbit formula; three characteristics & all selections \\
$\F_9$\newline $n=4$, $\lambda=1,i,-1$, $k=1$ & --- & $16,4,16$ & distributions from the orbit formula & all selections \\
$\F_9$, $\F_{25}$\newline $n=8$ and $n=6$, $\lambda=1$, $k=1$ & --- & $256,64$ & distributions from the orbit formula & all selections \\
$\F_{32}$\newline $n=31$, $\lambda=1$, $k=1$ & $(1,1)$, $(10,1)^3$ & $2^{31}$ & $248$ joint terms (Fig.~\ref{fig:joint}); mean $15/2$ & exact enumerator; individual codes \\
\bottomrule
\end{tabularx}
\end{table}

As an external benchmark, the Euclidean cyclic cases $(q,n)=(2,21)$ and
$(3,16)$ were computed with prime-field polynomial and matrix routines that do
not use the orbit machinery. Among all $64$ binary cyclic codes of length $21$
exactly $16$ have hull dimension $6$, and among all $128$ ternary cyclic codes of
length $16$ exactly $32$ have hull dimension $4$. These are the counts of
Examples~16 and~17 of \cite{sangwisut2015}, respectively, and the full joint
histograms of both families agree with the orbit enumeration.

\section{Conclusions}\label{sec:conclusion}

For a square-free affine algebra $\A$ with labeled field components and
simple-root moduli, every $\lambda$-constacyclic code over $\A$ corresponds to
exactly one labeled factor selection. For a compatible twist its $k$-Galois
hull dimension is a weighted boundary count on the reciprocal-factor orbits,
$\eta_k(C(J))=\sum_{s,O}w_O\sum_i\varepsilon_{O,i}(1-\varepsilon_{O,i+1})$ with
$w_O=m_sd_O$. The bijection needs only the square-free, simple-root and
fixed-label hypotheses.
The compatibility condition is needed for the same-factor-set permutation, hence
for the orbit product, but not for the dual-generator formula
(Theorem~\ref{thm:dual-generator}), and for incompatible components the hull is
trivial (Remark~\ref{rem:incompatible}).

Both statistics are additive over orbits, so the joint enumerator factorizes,
\[
\Epoly(u,z)=\prod_{s=1}^{N}\prod_{O\in\Os}
\operatorname{tr}\bigl(T_{w_O}(u,z)^{a_O}\bigr),
\]
This is more than the sum of a hull-only count and a dimension count. The two
statistics are carried by the same variables $u$ and $z$, so the coefficient
$[u^{\kappa}z^{\eta}]$ counts codes with both parameters prescribed.
Specialization recovers the code-dimension distribution, the hull polynomial,
the exact labeled LCD count $2^{\sum_s|\Os|}$, and the mean and variance of the
hull dimension, with the exceptional term for orbits of length two.
Differentiation in $z$ at fixed $\kappa$ gives the exact conditional
distribution, mean and variance of the hull dimension at every attainable code
dimension (Corollaries~\ref{cor:conditional}
and~\ref{cor:conditional-variance}).

The two descriptions of the one-orbit polynomial serve different purposes. The
transfer matrix $\operatorname{tr}(T_w^{a})$ gives a uniform recursion whose cost
grows polynomially in the number of factors and in the weight sum,
$O\bigl(N_F(W+1)(\lfloor W/2\rfloor+1)\bigr)$ integer additions
(Proposition~\ref{prop:algorithm}), whereas the closed form
$P_{a,w}(z)=2+\sum_r\frac{a}{r}\binom{a-1}{2r-1}z^{rw}$ exhibits each coefficient
directly as a count of run compositions. Both are exact and give the same
coefficients without enumerating the $2^{N_F}$ labeled selections.

The examples span the range of the formula. At one end,
$\F_4^{\times4}$ keeps all $4096$ selections within reach of hand computation,
and the mixed algebra $\F_4\times\F_{16}$ shows the role of the extension
weights. At the other end, $\F_{32}$ of length $31$ compresses the orbit data
$(1,1),(10,1)^3$ and $2^{31}$ selections into $248$ nonzero coefficients with
symmetric hull multiplicities (Fig.~\ref{fig:joint}). Direct finite-field linear-algebra computations reproduce
the orbit predictions in characteristics $2$, $3$ and $5$, including the
mixed-algebra case in quotient coordinates, and the published counts of
\cite{sangwisut2015} for the Euclidean cyclic cases $(q,n)=(2,21)$ and $(3,16)$
are recovered.

The framework has two limitations. First, it assumes simple-root moduli and
compatible twists; repeated-root settings \cite{dinh2010} are not covered, and
incompatible components, whose hull is trivial, lie outside the orbit product.
Second, codes are counted as labeled selections rather than as equivalence
classes; counting isometry classes would require a specified group action
\cite{vanLintWilson2001} and is a different problem. Within this scope the
transfer-matrix formulation suggests that further additive statistics of factor
selections, and the analogous enumerators over larger non-chain algebras, can be
handled in the same way.

\section*{Statements and Declarations}

\noindent\textbf{Funding.} No specific funding was received for this research.

\smallskip
\noindent\textbf{Competing interests.} The author declares no competing
interests.

\smallskip
\noindent\textbf{Author contributions.} Arun Kumar Yadav is the sole author,
designed the study, and reviewed and approved the manuscript.

\smallskip
\noindent\textbf{Data availability.} This work uses no empirical datasets. All
reported quantities are exact algebraic and combinatorial values determined by
the parameters listed in the article. The programs that generate the tables of
Section~\ref{sec:examples} and the generated coefficient tables are provided as
supplementary material.

\smallskip
\noindent\textbf{Use of AI-assisted tools.} An AI-assisted tool (Arena.ai Agent
Mode) was used to assist with mathematical checking, the development of the
verification code and language editing. The author reviewed and approved all
mathematical statements, computations and the final text, and takes full
responsibility for the content.

\appendix
\section{Reproducibility of the computations in Section~\ref{sec:examples}}\label{app:repro}

The orbit data are the inputs to the enumeration: the factorizations of the
$M_s$, the reciprocal permutation $\tau_{s,k}$, and the resulting orbit lengths
and degrees, from which the weights $w_O=m_sd_O$ are formed. All coefficients of
$\Epoly(u,z)$, the hull polynomials and the conditional moments are computed from
this data with exact integer and rational arithmetic, with no floating-point
arithmetic anywhere in the reported values.

Each family is evaluated in two independent ways: by the trace recurrence of
Proposition~\ref{prop:algorithm}, and through the closed-form orbit coefficients
of Proposition~\ref{prop:closed} followed by convolution. The two agree
coefficient by coefficient for every family of Table~\ref{tab:validation}. The
direct finite-field comparisons of Section~\ref{sec:verification} use
component-field linear algebra, as described there, together with the
quotient-coordinate treatment of the mixed algebra.

For $\F_{32}$ with $n=31$ the enumerator accounts for all $2^{31}$ selections
exactly through the orbit recurrence. Individual codes of that length were
compared with the orbit formula on a random subsample, which is an auxiliary
check of individual codes and not part of the count. The programs and the
coefficient tables they generate accompany the manuscript as supplementary
material and are not needed to state or verify the mathematics above.

\bibliographystyle{sn-mathphys-num-local}
\bibliography{references}

\end{document}
